% V5.1 patch applied 2026-03-25
% Fixes: B1 (abstract hedge), B2 (R3 Bonferroni), B5 (Liu 2024 softening),
%         C1 (expl2 vs think equivalence qualifier), C2 (Type I table cross-dataset label),
%         C4 (discordant pair clustering disclosure)
% Added: Appendix E (qualitative comparison), Appendix F (full per-model tables),
%        Appendix G (category breakdown), Appendix H (discordant pair analysis),
%        Iterative methodology section
% V5.2 patch applied 2026-03-25
% Added: difficulty-scaling hypothesis (abstract, intro, discussion subsection, conclusion)
%        Section 3.5 (representative trap examples: CC and DC)
%        Verbosity confound strengthened with token counts (limitation item 4)
%        Appendix I (DeepSeek median completion token counts)
\documentclass[11pt,a4paper]{article}

% -----------------------------------------------------------------------
%  Packages
% -----------------------------------------------------------------------
\usepackage[a4paper,margin=2.5cm]{geometry}
\usepackage{booktabs}
\usepackage{hyperref}
\usepackage{cleveref}
\usepackage[round,authoryear]{natbib}
\usepackage{graphicx}
\usepackage{enumitem}
\usepackage{amsmath}
\usepackage{microtype}
\usepackage{xcolor}
\usepackage{array}
\usepackage{multirow}
\usepackage{caption}

\hypersetup{
  colorlinks=true,
  linkcolor=blue!60!black,
  citecolor=blue!60!black,
  urlcolor=blue!60!black,
}

% -----------------------------------------------------------------------
%  Title block
% -----------------------------------------------------------------------
\title{\textbf{Forced Depth Consideration Reduces Type~II Errors\\
in LLM Self-Classification:\\
Evidence from an Exploration Prompting Ablation Study}}

\author{Wiktor Potapczyk\\
\small Independent Researcher\\
}

\date{March 2026}

% -----------------------------------------------------------------------
\begin{document}
\maketitle

% -----------------------------------------------------------------------
\begin{abstract}
LLM-based task classifiers systematically misclassify prompts that are
superficially simple but contextually complex, routing them to lightweight
processing paths (\emph{Type~II error}).  We introduce \emph{Step-0 prompt
framing}---prepending a single question before the classification
decision---as a novel intervention for LLM self-classification, and
present what is, to our knowledge, the first empirical investigation of
how this framing affects classification accuracy.

Using TaskClassBench (200 effective trap prompts, 4 models, 3 families),
we find that open-ended exploration (``What's really going on here?'')
reduces the pooled Type~II rate to 1.25\%, significantly outperforming
directed extraction at 3.12\% ($p < 0.001$, Bonferroni-corrected).
A mechanism ablation reveals that even a content-free metacognitive
instruction (``Think carefully'') achieves 1.0\%---not significantly
different from exploration ($p = 0.77$)---yet both significantly
outperform structured detection and extraction approaches.

The operative mechanism appears to be \emph{forced attention to task
complexity before classification}, rather than open-ended framing
specifically.  However, qualitative analysis of failure cases reveals
complementary failure modes: exploration forces committed implication
statements that anchor classification, while the metacognitive instruction
uses consequence framing that catches different trap subtypes.  Neither
approach has a universal advantage.

The effect is concentrated in models with moderate baseline error rates;
near-ceiling models show negligible benefit.  We hypothesise that this
capability gradient reverses at higher context loads ($>$100K tokens),
where even capable models would benefit from forced depth
consideration---a falsifiable prediction for future work.  All claims
are exploratory pending pre-registered replication.

\noindent\textbf{Keywords:} exploration prompting, task classification,
Type~II error, LLM routing, metacognitive prompting, depth consideration
\end{abstract}

% -----------------------------------------------------------------------
\section{Introduction}
\label{sec:intro}
% -----------------------------------------------------------------------

\subsection{Why task routing matters in LLM systems}
\label{sec:intro:routing}

Effective orchestration of large language model (LLM) systems increasingly
depends on accurate task classification: assigning incoming user requests
to appropriate processing paths, model tiers, or specialist agents.
\citet{Agrawal2025} demonstrate that task-level features predict routing
decisions, but that existing approaches rely on structured extraction of
surface features from the prompt text.  \citet{Huang2025} frame routing
as a forward-planning problem, showing that simulating downstream task
demands before committing to a route improves outcomes.  \citet{MarinelliPichlmeier2025}
propose using post-hoc chain-of-thought step counts (NofT) as a routing
signal.  \citet{ZhangRouterR12025} go further, training a reinforcement-learning
router that interleaves deliberation with routing decisions.

All these lines of work assume that the features driving routing decisions
are recoverable from the prompt at classification time.  We challenge this
assumption at the input level: \emph{how} the classifier is prompted to
consider those features may matter more than which features are targeted.

\subsection{The Type~II failure mode}
\label{sec:intro:type2}

We define \textbf{Type~II error} as the event in which a classifier assigns
a ``Quick'' (lightweight) label to a prompt that in fact requires deeper
processing: multi-step reasoning, contextual disambiguation, or expert-level
synthesis.  This is distinct from hallucination and from sycophancy,
though it may share a common mechanistic root.

\citet{Wu2024} show in the ProCo framework that classification models
learn spurious correlations from surface cues.  Our trap prompt categories
are constructed to exploit precisely this: prompts that look routine on
the surface but contain contextual depth signals.  \citet{MASFT2025}
document a reasoning-action mismatch in which models can recognise
complexity in chain-of-thought outputs but still act superficially.
Our work targets the classification decision \emph{before} any action
is taken.

\subsection{The exploration prompting hypothesis and its evolution}
\label{sec:intro:hypothesis}

\citet{Furniturewala2024} introduced \emph{Implication Prompting}: prepending
an implication question before answering a task
generates the bias rationale before requesting correction, matching
white-box debiasing methods on stereotypical response correction.  We extend
this to the meta-task of self-classification.  To our knowledge, no prior
work has investigated how Step-0 prompt framing---the question posed to the
model before it classifies its own task---affects the accuracy of LLM
self-classification decisions.  This gap was validated through systematic
search using Gemini~2.5~Pro (Deep Research) and Elicit.

Our investigation proceeded in four stages.  We began by comparing
open-ended exploration questions against directed intent extraction (Rounds
1 and 2), finding a significant advantage for open-ended framing once
confounds were controlled (EXP, $p = 0.000729$).  \citet{SycEval2025}
and \citet{Liu2024ICML} motivate this: directed step-by-step reasoning
can interfere with classification, while open-ended framing may break the
surface-agreement default.

We then ran a mechanism ablation (Round~3): if the advantage of open-ended
questions comes from forcing the model to \emph{engage with complexity}
rather than from the question format itself, a minimal content-free
instruction (``Think carefully about the complexity of this task'') should
produce not-significantly-different results.  R3 results are consistent
with this prediction (McNemar 5:7, $p = 0.774$).

This narrows the mechanism claim without weakening the practical
recommendation.  Structured extraction and detection remain significantly
worse.  R3 results suggest the advantage is broader than initially
framed: \emph{any instruction that forces unbounded engagement with task
complexity before classification achieves not-significantly-different
Type~II error rates on our benchmark}, though formal equivalence is not
established.

Our research question became: \emph{Does forced depth consideration---in
whatever form---reduce Type~II errors, and does open-ended framing confer
any additional advantage beyond the depth-engagement it triggers?}

This paper tests forced depth consideration under favourable conditions:
short prompts ($<$512 tokens), minimal prior context, and moderate
classification difficulty.  Under these conditions, a content-free
metacognitive directive achieves near-equivalent performance to open-ended
exploration (Section~6).  However, the per-model results reveal a
capability gradient that motivates a broader hypothesis: if forced depth
consideration acts as an external scaffold for the model's implicit
complexity assessment, then the intervention's value should scale with the
difficulty of that assessment.  As context windows grow and classification
inputs become embedded in extensive conversational history, even capable
models may face complexity-assessment demands that exceed their implicit
capacity.  This scaling prediction is not tested here; it is stated as a
falsifiable hypothesis that structures our interpretation.

% -----------------------------------------------------------------------
\section{Background and Related Work}
\label{sec:background}
% -----------------------------------------------------------------------

\subsection{Implication prompting}

\citet{Furniturewala2024} showed that prepending an implication question to
QA tasks generates the bias rationale before requesting correction,
matching white-box debiasing methods on stereotypical response correction.
Our contribution applies a structurally similar strategy to the meta-task
of self-classification: the model is classifying the \emph{nature of the
request}, not answering the request itself.

\subsection{Lookahead and task-aware routing}

\citet{Huang2025} propose lookahead routing, in which a model simulates the
demands of a task before selecting a processing path.  We operationalise a
lightweight variant: one question (or a metacognitive directive) replaces
full forward simulation.  Our Round~3 result suggests that even a
content-free nudge can approximate the benefit.

\citet{Agrawal2025} establish task-level features as routing signals.
Our results suggest that \emph{how those features are elicited} may matter
as much as which features are targeted: the same model, given the same
prompt, surfaces different features depending on Step-0 framing.

\citet{MarinelliPichlmeier2025} use post-hoc NofT metadata to route after
generation.  Our approach acts \emph{before} generation via prompt framing,
making it applicable to single-turn classifiers that never produce
long-form CoT.  \citet{ZhangRouterR12025} integrate deliberation into the
router itself via reinforcement learning; our ablation result suggests
that even an un-trained metacognitive directive achieves substantial
improvement.

\subsection{Chain-of-thought and classification}

A substantial literature documents that CoT prompting improves reasoning
tasks \citep{Wei2022,Kojima2022}.  However, \citet{Liu2024ICML} showed that
CoT can \emph{reduce} classification performance by up to 36\% on tasks
where implicit pattern-matching outperforms deliberation.  This finding is
consistent with the verbal-overshadowing effect in human cognition
\citep{Liu2024ICML}, providing a conceptual parallel from human cognition
research.  We note that \citet{Liu2024ICML} studied human verbal
overshadowing, not LLM binary classification; the analogy is suggestive,
not direct evidence.  Our gate sensitivity pattern is nonetheless consistent
with this account: structured yes/no depth detection forces the model to
deliberate in a constrained format that may interfere with the pattern-based
recognition required for trap prompts.

\citet{GunasekaraRatnayake2025} found that unstructured natural-language
reasoning outperforms structured JSON reasoning plans by up to 19\%
on the MATH benchmark (zero-shot unstructured outperforming five-shot
structured).  Our open vs.\ directed finding is directionally consistent:
constraining the output format of the Step-0 response appears to reduce
its value.

\subsection{Capability interaction and prompt sensitivity}

\citet{Khan2025} documents a ``Prompting Inversion'': constraints that help
mid-tier models can hurt frontier models by inducing hyper-literalism.
This contextualises our gate sensitivity finding---structured depth
detection helps DeepSeek but harms Claude models---and provides a
framework for the hypothesis that exploration and metacognitive advantages
may be capability-moderated (\Cref{sec:future}).

\subsection{Classification failure modes}

\citet{Wu2024} demonstrate that spurious surface correlations systematically
bias classification.  Our trap prompt categories are constructed to create a
mismatch between surface appearance and required depth, making them a direct
empirical test of the surface-correlation failure mode.

\citet{MASFT2025} report a reasoning-action mismatch: models can produce
chain-of-thought that correctly identifies task complexity but still
output a simple action label.  Our gated variant (``Are depth signals
present? yes/no'') tests this directly: the structured output format may
amplify the mismatch by forcing a binary answer.

\subsection{Sycophancy and surface-default behaviour}

\citet{SycEval2025} document that LLMs exhibit persistent sycophantic
responses to surface-level task framing.  Our baseline condition captures
this failure operationally: prompts that look simple on the surface are
systematically assigned to the ``Quick'' category even when their contextual
depth warrants deeper processing.  The open-ended pre-question and the
metacognitive directive are both candidate interventions.

% -----------------------------------------------------------------------
\section{Task Setup and Benchmark Design}
\label{sec:benchmark}
% -----------------------------------------------------------------------

\subsection{Experimental rounds overview}

The study proceeded in four independent API rounds (\Cref{tab:rounds}).
All rounds used temperature~0; \Cref{sec:stability} quantifies the
residual run-to-run variance.

\begin{table}[htbp]
\centering
\caption{Experimental rounds.  R1, R2, EXP, and R3 are separate API runs.
  Temperature-0 non-determinism produces 1--2 error differences across
  runs (verified; \Cref{sec:stability}).  Numbers across datasets
  may not be combined.
  [\emph{Source: findings-structured.md, data collection section}]}
\label{tab:rounds}
\smallskip
\small
\resizebox{\textwidth}{!}{%
\begin{tabular}{@{}llrrll@{}}
\toprule
Round & Dataset & Models & Traps ($N$) & Variants & Key design note \\
\midrule
R1 (pilot)     & R1  & 5 (incl.\ Gemini Pro) & 120 & base, expl, extr
  & \emph{Confounded}: unequal Quick gates \\
R2 (controlled) & R2 & 5 (incl.\ Gemini Pro) & 120 & gated, extr2, expl2, expl3
  & Equal gates; resolves R1 confound \\
EXP (primary)  & EXP & 4 (excl.\ Gemini Pro) & 200 eff.\ & all 7 variants
  & Post-hoc expansion after R2 ($p=0.065$) \\
R3 (mechanism ablation) & R3 & 4 (same as EXP) & 200 eff.\ & think (1 variant)
  & Cross-dataset comparison with EXP \\
\bottomrule
\end{tabular}}% end resizebox
\end{table}

\subsection{Iterative experimental design}
\label{sec:benchmark:iterative}

This study followed an iterative, self-correcting design process rather than
a single pre-registered protocol.  We document this explicitly because the
design choices at each stage were driven by findings from the previous stage,
and because transparency about this process is required for correct
interpretation of statistical claims.

\begin{enumerate}[itemsep=2pt]

  \item \textbf{R1 (Pilot).}  R1 was designed to test whether open-ended
    exploration questions reduce Type~II errors.  It discovered a strong
    exploration advantage (expl vs.\ base, $p = 0.000977$; expl vs.\ extr,
    $p = 0.000002$), but post-hoc analysis by a panel of 5 LLM reviewers
    identified three confounds: (a) the Quick gate was looser for
    \textsf{base} than for \textsf{expl}/\textsf{extr}, confounding gate
    tightness with question framing; (b) the \textsf{extr} control asked the
    model to name a specialist agent---a premature routing step rather than
    intent extraction; (c) no within-family wording ablation was included.

  \item \textbf{R2 (Controlled).}  R2 equalized the Quick gate across all
    variants and replaced the poor extraction control (\textsf{extr}) with a
    cleaner one (\textsf{extr2}: ``Summarize the user's intent in one
    sentence'').  Three open-ended wordings were tested (\textsf{expl2},
    \textsf{expl3}, plus retained \textsf{expl}).  R2 showed a directional
    trend ($p = 0.065$ for \textsf{expl2} vs.\ \textsf{extr2}, 4-model) but
    insufficient power at $N = 120$ trap prompts to support a positive claim.
    Gate sensitivity (Claude~Haiku and Claude~Sonnet harmed by structured
    binary detection) emerged as a secondary finding.

  \item \textbf{EXP (Power expansion).}  EXP expanded the benchmark
    post-hoc by increasing CC and DC categories from 20 to 100 prompts each
    (200 effective traps per model), motivated by the R2 power shortfall.
    This expansion was not pre-registered; the choice of which categories to
    expand was informed by R1 and R2 discrimination patterns.  EXP is
    therefore an exploratory study, not a confirmatory one.

  \item \textbf{R3 (Mechanism ablation).}  Following EXP's finding that
    \textsf{expl2} outperforms \textsf{extr2}, the competing hypothesis
    (``any instruction that forces depth consideration achieves the same
    result'') was addressed by a dedicated mechanism ablation.  R3 was
    motivated by adversarial review of the EXP finding, which noted that the
    expl2 advantage could be explained by forced engagement with complexity
    rather than by open-ended framing per se.

\end{enumerate}

\paragraph{Why iterative design is a feature, not a limitation.}
Post-hoc expansion and adversarial re-analysis are often treated as
methodological weaknesses.  We argue they are epistemically responsible:
each round explicitly tests the strongest available critique of the previous
round's finding, and each design choice is documented with its motivating
analysis.  The alternative---treating R1 as confirmatory and ignoring the
confounds identified by review---would have produced a cleaner-looking result
with a weaker epistemological foundation.  The cost is that all claims must
be treated as exploratory pending pre-registered replication; this cost is
stated explicitly throughout this paper.

\subsection{Variant taxonomy}

\Cref{tab:variants} lists the seven EXP prompt variants and the R3 mechanism
variant.  Each variant adds one Step-0 sentence before the classification
instruction.  The baseline (\textsf{base}) adds no Step-0 question.

\begin{table}[htbp]
\centering
\caption{Variant taxonomy.
  [\emph{Source: findings-structured.md, variant design section}]}
\label{tab:variants}
\smallskip
\resizebox{\textwidth}{!}{%
\begin{tabular}{@{}llp{5.2cm}ll@{}}
\toprule
Variant & Round & Step-0 text & Gate & Category \\
\midrule
\textsf{base}   & R1 & None & Loose (R1 confound) & Baseline \\
\textsf{expl}   & R1 & ``What does this prompt imply?'' & Tighter & Open \\
\textsf{extr}   & R1 & ``Name the best specialist agent'' & Tighter
  & Directed (poor control) \\
\textsf{gated}  & R2 & ``Are depth signals present? yes/no'' & Equal & Structured \\
\textsf{extr2}  & R2 & ``Summarize the user's intent in one sentence'' & Equal
  & Directed (primary control) \\
\textsf{expl2}  & R2 & ``What's really going on here?'' & Equal
  & \textbf{Open (primary treatment)} \\
\textsf{expl3}  & R2 & ``What does this request actually require?'' & Equal
  & Open (wording variant) \\
\textsf{think}  & R3 & ``Think carefully about the complexity of this task'' & Equal
  & \textbf{Metacognitive (mechanism ablation)} \\
\bottomrule
\end{tabular}}% end resizebox
\end{table}

\noindent The R1 confound arises because \textsf{base} uses a looser Quick
gate than \textsf{expl} and \textsf{extr}.  R2 resolves this by assigning
identical gates to all variants.  The R2 \textsf{extr} control is also
improved: ``Name the best specialist agent'' forces premature routing rather
than intent extraction; \textsf{extr2} is a cleaner extraction control.

\subsection{Benchmark categories}

\Cref{tab:categories} describes the six trap categories and one control
category.  The two \emph{discriminative} categories (context-contradiction,
CC; disguised-correction, DC) are the only ones that produced meaningful
variation across variants in all datasets.  These 200 effective prompts
(100 CC + 100 DC per model) form the primary analysis set.
The four \emph{non-discriminative} categories produced near-zero errors
across all variants in all datasets; they are retained as completeness
checks but excluded from the primary Type~II analysis.

\begin{table}[htbp]
\centering
\caption{Benchmark categories in EXP.
  [\emph{Source: EXP metadata, calculations-expanded.json, 4 models,
  $N_\text{eff} = 200$}]}
\label{tab:categories}
\smallskip
\resizebox{\textwidth}{!}{%
\begin{tabular}{@{}lrll@{}}
\toprule
Category & $N$ per model (EXP) & Role & Notes \\
\midrule
context-contradiction (CC) & 100 & Discriminative & Expanded from 20 in R1/R2 \\
disguised-correction (DC)  & 100 & Discriminative & Expanded from 20 in R1/R2 \\
cross-reference            &  20 & Non-discriminative & Near-zero errors all variants \\
hidden-compound            &  20 & Non-discriminative & Near-zero errors all variants \\
perspective-shift          &  20 & Non-discriminative & Near-zero errors all variants \\
vague-affect               &  20 & Non-discriminative & Near-zero errors all variants \\
genuinely-quick (control)  &  31 & Type~I control & Not counted in Type~II analysis \\
\midrule
\multicolumn{2}{@{}l}{\textbf{Effective trap set (CC + DC)}} & \textbf{200} & Primary analysis \\
\multicolumn{2}{@{}l}{All-trap set (incl.\ non-discriminative)} & 280
  & Completeness check only \\
\bottomrule
\end{tabular}}
\end{table}

\paragraph{Post-hoc expansion disclosure.}
EXP was created after R2 yielded $p = 0.065$ (insufficient power at
$N = 120$ trap prompts per model).  We expanded CC and DC from 20 to 100
prompts each to increase statistical power.  This expansion was motivated
by the R2 result, not pre-registered.  All EXP results must therefore be
treated as exploratory pending pre-registered replication with an
independent sample.

\subsection{Representative trap examples}
\label{sec:benchmark:examples}

The following examples illustrate why the two discriminative categories
(context-contradiction and disguised-correction) are the primary sources
of Type~II error.  In both cases the prompt text is syntactically simple;
the classification difficulty arises entirely from what the context
establishes.

\subsubsection*{Example 1: context-contradiction (CC)}

\noindent\textbf{Context:}
\begin{quote}
\textit{``You are building an enterprise SaaS dashboard with real-time
analytics.  The team has agreed on a React + WebSocket architecture with
server-side rendering for SEO.''}
\end{quote}

\noindent\textbf{Prompt:}
\begin{quote}
\textit{``Actually, let's just use a static site generator.''}
\end{quote}

\noindent\textbf{Why it is a trap.}
The prompt reads as a simple architectural preference.  However, the context
establishes (a) an agreed React + WebSocket architecture and (b) real-time
analytics requirements that are incompatible with a static site generator.
Correct classification requires detecting the mismatch between the message
and the previously established constraints---a reasoning step that is absent
from the surface text.  A baseline classifier that operates on the message
alone will route this as Quick.

\subsubsection*{Example 2: disguised-correction (DC)}

\noindent\textbf{Context:}
\begin{quote}
\textit{``You are developing a fault-tolerant ETL pipeline with retry logic,
dead-letter queues, and automated alerting.''}
\end{quote}

\noindent\textbf{Prompt:}
\begin{quote}
\textit{``we don't need the retry logic actually''}
\end{quote}

\noindent\textbf{Why it is a trap.}
Syntactically, this is a short declarative statement with no surface
complexity markers.  Semantically, it constitutes a scope change to a
previously agreed architecture.  Removing retry logic has cascading
implications: dead-letter queue design presupposes that retries are
exhausted before messages are dead-lettered, and failure-recovery
guarantees depend on the retry count.  The user is revising a design
commitment, not asking a simple question.  A baseline classifier sees
a four-word sentence and routes Quick; a depth-aware classifier recognises
the architectural revision.

% -----------------------------------------------------------------------
\section{Experimental Evolution: Pilot and Controlled Rounds}
\label{sec:evolution}
% -----------------------------------------------------------------------

This section presents R1 and R2 as historical context establishing why
the primary expanded dataset (EXP) was necessary.  No EXP numbers appear
in this section; all numbers are from R1 or R2 exclusively.

\subsection{R1: Pilot (confounded gates)}

\Cref{tab:r1} shows per-model Type~II error counts for R1.
\textsf{expl} produced dramatically fewer errors than both \textsf{base}
and \textsf{extr} in every model where errors were present.

\begin{table}[htbp]
\centering
\caption{R1 per-model Type~II error counts (confounded pilot).
  \textbf{Caution}: \textsf{base} uses a looser Quick gate than
  \textsf{expl}/\textsf{extr}; the expl--base comparison is
  \emph{not} a clean test of prompt framing.
  [\emph{Source: R1 = calculations-r1.json, 5 models,
  $N = 120$ trap prompts per model}]}
\label{tab:r1}
\smallskip
\begin{tabular}{@{}lrrr@{}}
\toprule
Model & \textsf{base} (/120) & \textsf{expl} (/120) & \textsf{extr} (/120) \\
\midrule
DeepSeek       &  7 & \textbf{0}  & 15 \\
Gemini Flash   &  2 & \textbf{0}  &  2 \\
Gemini Pro     &  0 &             0  &  1 \\
Claude Haiku   &  5 & \textbf{1}  &  4 \\
Claude Sonnet  &  4 &             4  &  6 \\
\midrule
Pooled (5 models) & 18 (3.0\%) & \textbf{5 (0.83\%)} & 28 (4.67\%) \\
\bottomrule
\end{tabular}
\end{table}

\noindent Pooled McNemar tests (5 models, 120 traps each):
\begin{itemize}[noitemsep]
  \item \textsf{expl} vs.\ \textsf{base}: 14:1 discordant,
    $p = 0.000977$ [\emph{R1, mcnemar\_tests.expl\_vs\_base}]
  \item \textsf{expl} vs.\ \textsf{extr}: 24:1 discordant,
    $p = 0.000002$ [\emph{R1, mcnemar\_tests.expl\_vs\_extr}]
\end{itemize}

Despite the strong signal, R1 results cannot be taken as clean evidence for
prompt framing.  The confound (unequal gates) means the observed advantage
of \textsf{expl} over \textsf{base} conflates gate tightness with question
framing.  Additionally, \textsf{extr} forces premature routing by asking the
model to name a specialist agent---a poor extraction control that may itself
introduce errors.  R1 motivates R2; it does not support the primary claim.

\subsection{R2: Controlled round (equal gates)}

R2 equalized the Quick gate across all variants, yielding a clean comparison
between framing strategies (\Cref{tab:r2}).  Gate sensitivity emerged as a
secondary finding: Claude~Haiku gated errors (17) far exceeded its \textsf{expl2}
errors (0), and similarly for Claude~Sonnet (15 vs.\ 3).

\begin{table}[htbp]
\centering
\caption{R2 per-model Type~II error counts (controlled, equal gates).
  [\emph{Source: R2 = calculations-r2.json, 5 models,
  $N = 120$ trap prompts per model, equal Quick gates}]}
\label{tab:r2}
\smallskip
\begin{tabular}{@{}lrrrr@{}}
\toprule
Model & \textsf{gated} (/120) & \textsf{extr2} (/120)
      & \textsf{expl2} (/120) & \textsf{expl3} (/120) \\
\midrule
DeepSeek       &  3 &  2 & \textbf{1}  & \textbf{0} \\
Gemini Flash   &  1 &  1 & \textbf{0}  &  1 \\
Gemini Pro     &  1 & \textbf{0} &  1 & \textbf{0} \\
Claude Haiku   & 17 &  4 & \textbf{0}  &  2 \\
Claude Sonnet  & 15 &  4 &  3          & \textbf{1} \\
\midrule
Pooled 5-model & 37 (6.17\%) & 11 (1.83\%) & 5 (0.83\%) & 4 (0.67\%) \\
Pooled 4-model (excl.\ Pro) & 36 (7.5\%) & 11 (2.29\%) & 4 (0.83\%) & 4 (0.83\%) \\
\bottomrule
\end{tabular}
\end{table}

\noindent McNemar tests for the primary comparison (4-model, excluding Gemini Pro):
\begin{itemize}[noitemsep]
  \item \textsf{expl2} vs.\ \textsf{extr2}: 9:2 discordant,
    $p = 0.065$ --- \textbf{not significant}
    [\emph{R2, mcnemar\_4model.expl2\_vs\_extr2}]
  \item \textsf{expl2} vs.\ \textsf{gated}: 32:0 discordant,
    $p < 0.001$ [\emph{R2, mcnemar\_4model.expl2\_vs\_gated}]
\end{itemize}

Including Gemini Pro (5-model pooling), the primary comparison is
$p = 0.146$ [\emph{R2, mcnemar\_5model.expl2\_vs\_extr2}].

R2 isolates prompt framing as the operative variable and shows exploration
trending better than extraction, but $p = 0.065$ is insufficient to support
a positive claim at $N = 120$.  The gate sensitivity pattern (Haiku and
Sonnet harmed by the structured gate) emerges as a secondary phenomenon
warranting investigation.  These two findings together motivated the power
expansion in EXP.

% -----------------------------------------------------------------------
\section{Primary Analysis: Expanded Benchmark}
\label{sec:primary}
% -----------------------------------------------------------------------

EXP is the primary evidence.  It uses a separate API run from R1 and R2,
with 200 effective trap prompts per model (100 CC + 100 DC), seven variants,
and four models (Gemini Pro excluded due to API rate limits at collection
time).  All numbers in this section are from \texttt{calculations-expanded.json}.

\subsection{Per-model error counts}

\Cref{tab:exp_errors} shows effective Type~II errors (out of 200 effective
traps per model) for all seven variants.

\begin{table}[htbp]
\centering
\caption{EXP per-model effective Type~II error counts (primary evidence).
  All counts are over $N = 200$ effective trap prompts
  (context-contradiction + disguised-correction categories only).
  \textsf{expl2} column is the primary treatment; \textsf{extr2} is the
  primary control.
  [\emph{Source: EXP = calculations-expanded.json, error\_counts section,
  4 models, $N_\text{eff} = 200$}]}
\label{tab:exp_errors}
\smallskip
\resizebox{\textwidth}{!}{%
\begin{tabular}{@{}lrrrrrrr@{}}
\toprule
Model
  & \textsf{base}
  & \textsf{expl}
  & \textsf{extr}
  & \textsf{gated}
  & \textsf{extr2}
  & \textsf{expl2}
  & \textsf{expl3} \\
\midrule
DeepSeek       & 16 &  2 & 67 &  5 &  8 & \textbf{2}  & \textbf{0} \\
Gemini Flash   &  3 &  0 &  5 &  2 &  1 & \textbf{0}  &  3 \\
Claude Haiku   & 10 &  2 &  5 & 43 &  8 & \textbf{1}  &  5 \\
Claude Sonnet  & 12 & 12 & 16 & 34 &  8 & \textbf{7}  & 11 \\
\midrule
\textbf{Pooled (4 models)}
  & 41 (5.12\%) & 16 (2.00\%) & 93 (11.62\%) & 84 (10.50\%)
  & 25 (3.12\%) & \textbf{10 (1.25\%)} & 19 (2.38\%) \\
\bottomrule
\end{tabular}}% end resizebox
\end{table}

\noindent A supporting result: \textsf{expl2} also outperforms the no-question baseline
(\textsf{base}): 32:1 discordant pairs, $p < 0.001$, surviving Bonferroni correction
[\emph{EXP, mcnemar\_tests.expl2\_vs\_base}].

\paragraph{Discordant pair clustering (C4 disclosure).}
The 17 expl2-wins discordant pairs come from 14 unique prompt indices across 4
models; 3 prompts contribute discordant pairs in 2 models (prompts \#5, \#11,
\#194), and 11 prompts contribute in only 1 model.  The effective independent
observation count is thus closer to 14 unique prompts than 17 individual
pairs.  The McNemar test operates on matched pairs within each model stratum
(as confirmed by the CMH stratification; \Cref{tab:cmh}), which accounts
for the model-level pairing structure.  The 2 extr2-wins pairs come from
2 unique prompts, both from Claude~Sonnet (\#5 and \#22).  Notably, prompt
\#5 contributes expl2-wins on DeepSeek and Claude~Haiku but an extr2-win on
Claude~Sonnet, illustrating cross-model heterogeneity.  Full discordant pair
detail is provided in \Cref{app:discordant}.

\noindent Notable observations from \Cref{tab:exp_errors}:
\begin{itemize}[noitemsep]
  \item \textsf{expl2} achieves the lowest pooled error
    rate (1.25\%) among the seven EXP variants.
  \item \textsf{extr} (the R1-era directed control) catastrophically
    increases DeepSeek errors from 16 to 67---confirming that
    ``Name the best specialist agent'' is a poor extraction control.
  \item \textsf{gated} catastrophically increases Claude~Haiku errors
    from 10 to 43 and Claude~Sonnet errors from 12 to 34; this is the
    gate sensitivity finding (\Cref{sec:gate}).
  \item Claude~Sonnet shows the weakest response to \textsf{expl2}
    (7 errors vs.\ 8 for \textsf{extr2}), consistent with its
    individually non-significant per-model McNemar result.
\end{itemize}

\subsection{Hypothesis tests}

\Cref{tab:mcnemar} presents the primary McNemar tests and their Bonferroni
status.  We tested 8 comparisons against the EXP dataset; the Bonferroni
threshold is $\alpha = 0.05 / 8 = 0.00625$.

\begin{table}[htbp]
\centering
\caption{Primary McNemar tests, EXP.  All comparisons are pooled across
  4 models and 200 effective trap prompts per model (800 total observations).
  Bonferroni correction: 8 comparisons, $\alpha = 0.00625$.
  [\emph{Source: EXP = calculations-expanded.json,
  mcnemar\_tests and bonferroni sections, 4 models,
  $N_\text{eff} = 200$}]}
\label{tab:mcnemar}
\smallskip
\begin{tabular}{@{}lrrrll@{}}
\toprule
Comparison & \textsf{a} wins & \textsf{b} wins & $p$-value
  & Bonferroni & Note \\
\midrule
\textsf{expl2} vs.\ \textsf{extr2} (CORE)
  & 17 & 2 & 0.000729 & \textbf{YES} & Primary claim \\
\textsf{expl2} vs.\ \textsf{gated}
  & 74 & 0 & $<$0.001  & \textbf{YES} & Gate sensitivity \\
\textsf{expl2} vs.\ \textsf{base}
  & 32 & 1 & $<$0.001  & \textbf{YES} & vs.\ no question \\
\textsf{extr2} vs.\ \textsf{gated}
  & 68 & 9 & $<$0.001  & \textbf{YES} & Extraction $>$ detection \\
\textsf{expl2} vs.\ \textsf{expl3} (wording)
  & 14 & 5 & 0.064 & no & Hypothesis only \\
\textsf{expl} vs.\ \textsf{expl2} (wording)
  &  1 & 7 & 0.070 & no & Hypothesis only \\
\textsf{expl} vs.\ \textsf{expl3} (wording)
  &  9 & 6 & 0.607 & no & Equivalent \\
\bottomrule
\end{tabular}
\end{table}

\subsection{CMH stratified analysis: core claim}

The naive pooled McNemar treats all 800 prompt-model pairs as exchangeable.
The Cochran--Mantel--Haenszel (CMH) test stratifies by model, accounting
for the heterogeneous baseline error rates.  CMH is the more conservative
test; both agree (\Cref{tab:cmh}).

\begin{table}[htbp]
\centering
\caption{CMH stratified analysis for the core comparison
  (\textsf{expl2} vs.\ \textsf{extr2}), stratified by model.
  $cw$ = concordant win for \textsf{expl2} (expl2 correct,
  extr2 incorrect); $wc$ = concordant win for \textsf{extr2}.
  [\emph{Source: EXP = calculations-expanded.json,
  cmh\_tests.expl2\_vs\_extr2, 4 strata}]}
\label{tab:cmh}
\smallskip
\begin{tabular}{@{}lrr@{}}
\toprule
Model (stratum) & \textsf{expl2} wins ($cw$) & \textsf{extr2} wins ($wc$) \\
\midrule
DeepSeek       & 6 & 0 \\
Gemini Flash   & 1 & 0 \\
Claude Haiku   & 7 & 0 \\
Claude Sonnet  & 3 & 2 \\
\midrule
\textbf{CMH} $\chi^2$        & \multicolumn{2}{r}{10.32} \\
\textbf{CMH} $p$-value       & \multicolumn{2}{r}{0.0013} \\
Common odds ratio             & \multicolumn{2}{r}{8.5} \\
Breslow--Day $p$-value        & \multicolumn{2}{r}{0.0997} \\
\bottomrule
\end{tabular}
\end{table}

\noindent The Breslow--Day statistic tests homogeneity of the odds ratio
across strata.  $p = 0.0997$ does not reject homogeneity at $\alpha = 0.05$,
indicating that the effect direction is consistent across models.  However,
the $p$-value is close to the threshold and the per-stratum counts are small,
so the homogeneity result should be interpreted cautiously.

\subsection{Per-model McNemar: core comparison}

\Cref{tab:permodel} shows the per-model McNemar for \textsf{expl2} vs.\
\textsf{extr2}.  The pooled result is driven by DeepSeek and Claude~Haiku;
Gemini~Flash and Claude~Sonnet are individually non-significant.

\begin{table}[htbp]
\centering
\caption{Per-model McNemar tests, \textsf{expl2} vs.\ \textsf{extr2} (EXP).
  [\emph{Source: EXP = calculations-expanded.json,
  per\_model\_mcnemar section, $N_\text{eff} = 200$ per model}]}
\label{tab:permodel}
\smallskip
\resizebox{\textwidth}{!}{
\begin{tabular}{@{}lrrrll@{}}
\toprule
Model & \textsf{expl2} wins & \textsf{extr2} wins & $p$-value
  & Significant? & Note \\
\midrule
DeepSeek      & 6 & 0 & 0.031   & YES\,* & \\
Gemini Flash  & 1 & 0 & 1.0     & no     & Near-ceiling at baseline \\
Claude Haiku  & 7 & 0 & 0.016   & YES\,* & \\
Claude Sonnet & 3 & 2 & 1.0     & no     & Mixed signal \\
\bottomrule
\multicolumn{6}{@{}l}{\footnotesize
  * Individual significance at $p < 0.05$; note that per-model tests
  are underpowered and these $p$-values are not Bonferroni-corrected.}
\end{tabular}}
\end{table}

This heterogeneity is an important qualifier on the primary claim: the
statement ``exploration prompting reduces Type~II errors'' is better scoped
as ``exploration prompting reduces Type~II errors \emph{for models with
moderate baseline error rates}.''  Gemini~Flash approaches ceiling at
baseline (3 errors / 200), leaving little room for improvement.
Claude~Sonnet shows a mixed signal (3:2 discordant pairs) with the pooled
count directionally consistent but individually non-significant.

\subsection{Robustness check (supplementary)}

Excluding Claude~Sonnet (the model with the weakest and most mixed signal),
the core comparison becomes 14:0 discordant pairs, $p = 0.000122$
[\emph{EXP, mcnemar\_excl\_sonnet.expl2\_vs\_extr2, 3 models}].
We present this as a supplementary result only.  Excluding the model that
attenuates the effect is a selective robustness check; it demonstrates
that the result is not fragile, but it cannot substitute for the
full-sample analysis.

% -----------------------------------------------------------------------
\section{Round 3: Mechanism Ablation}
\label{sec:r3}
% -----------------------------------------------------------------------

The EXP result established that open-ended exploration outperforms
directed extraction.  The competing explanation---that \emph{any} instruction
forcing the model to consider task complexity achieves the same result,
regardless of question framing---was addressed by a dedicated mechanism
ablation.  \textbf{R3 is a cross-dataset comparison with EXP}: same prompts
(\texttt{prompts-expanded.yaml}), same four models, but a separate API run
conducted independently.  Cross-dataset comparisons carry 1--2 error
variation risk from temperature-0 non-determinism (\Cref{sec:stability}).

\subsection{Think-carefully variant}

The \textsf{think} variant prepends: ``Think carefully about the complexity
of this task.''  This is a content-free metacognitive directive.  It imposes
no specific question to answer, no format to fill, and no reasoning pathway
to follow.  It simply instructs the model to consider complexity before
classifying.

\subsection{Per-model results}

\Cref{tab:r3_errors} shows effective Type~II errors for \textsf{think}
across all four models.

\begin{table}[htbp]
\centering
\caption{R3 per-model effective Type~II error counts (\textsf{think} variant).
  \textbf{Note}: R3 is a separate API run from EXP; these figures are not
  directly comparable to EXP within-dataset figures.  Cross-dataset
  comparison carries $\pm$1--2 error variance.
  [\emph{Source: R3 = calculations-r3.json, error\_counts section,
  4 models, $N_\text{eff} = 200$ per model}]}
\label{tab:r3_errors}
\smallskip
\begin{tabular}{@{}lrrrr@{}}
\toprule
Model & Type~II errors (/200 eff.) & Type~II \% & Type~I errors (/31) & Type~I \% \\
\midrule
DeepSeek       & 2 & 1.00 &  7 & 22.58 \\
Gemini Flash   & 1 & 0.50 &  6 & 19.35 \\
Claude Haiku   & 1 & 0.50 &  6 & 19.35 \\
Claude Sonnet  & 4 & 2.00 &  2 &  6.45 \\
\midrule
\textbf{Pooled} & \textbf{8/800 (1.00\%)} & & 21/124 & 16.94\% \\
\bottomrule
\end{tabular}
\end{table}

\noindent Wilson 95\% CIs for \textsf{think} pooled: point 1.0\%,
computed from per-model data [\emph{R3, wilson\_cis section}].

\subsection{McNemar comparisons (cross-dataset)}

\textbf{Critical caveat}: all comparisons in \Cref{tab:r3_mcnemar} are
\emph{cross-dataset}---EXP run versus R3 run.  The McNemar test is applied
across matched prompts but different API calls.  Temperature-0
non-determinism adds up to 2 error differences; this affects the
smallest counts but does not explain the magnitude of the significant
comparisons.

\begin{table}[htbp]
\centering
\caption{Cross-dataset McNemar comparisons: EXP variants vs.\ R3
  \textsf{think}.  All 4 models, pooled across 200 effective traps per
  model (800 total observations).
  \textbf{Cross-dataset}: EXP run vs.\ R3 run, same prompts.
  [\emph{Source: R3 = calculations-r3.json, mcnemar\_vs\_exp section}]}
\label{tab:r3_mcnemar}
\smallskip
\resizebox{\textwidth}{!}{
\begin{tabular}{@{}llrrrll@{}}
\toprule
Comparison & Direction & EXP wins & R3 wins & $p$-value & Significant? & Interpretation \\
\midrule
\textsf{expl2} vs.\ \textsf{think}
  & --- & 5 & 7 & 0.774 & no & \textbf{Not sig.\ (see note)} \\
\textsf{expl3} vs.\ \textsf{think}
  & think better & 4 & 15 & 0.019 & YES & think $>$ expl3 \\
\textsf{extr2} vs.\ \textsf{think}
  & think better & 4 & 21 & 0.0009 & YES & think $\gg$ extr2 \\
\textsf{base} vs.\ \textsf{think}
  & think better & 1 & 34 & $<$0.001 & YES & think $\gg$ base \\
\textsf{gated} vs.\ \textsf{think}
  & think better & 1 & 77 & $<$0.001 & YES & think $\gg$ gated \\
\bottomrule
\end{tabular}}
\end{table}

\noindent Key results:
\begin{itemize}[noitemsep]
  \item \textsf{expl2} ($10/800$, $1.25\%$) and \textsf{think} ($8/800$,
    $1.00\%$) are not significantly different (McNemar 5:7, $p = 0.774$,
    12 discordant pairs).  With only 12 discordant pairs, a formal
    equivalence test (TOST) lacks power; the non-significant result reflects
    both possible equivalence and insufficient sample size to detect a small
    effect.  We do not claim formal statistical equivalence.
  \item \textsf{think} significantly outperforms \textsf{extr2}
    (4:21 discordant, $p = 0.0009$) and \textsf{expl3}
    (4:15 discordant, $p = 0.019$).
  \item On Claude~Sonnet, \textsf{think} (4 errors) performs numerically
    better than \textsf{expl2} (7 errors in EXP); the per-model McNemar
    is 3:6 discordant, $p = 0.508$ (not significant at individual level).
  \item The Type~I rate for \textsf{think} (16.94\%) exceeds
    \textsf{expl2}'s (12.90\%), discussed in \Cref{sec:typei}.
\end{itemize}

\subsection{Bonferroni correction within R3}
\label{sec:r3:bonferroni}

R3 involves 5 cross-dataset comparisons (against the EXP variants
\textsf{expl2}, \textsf{expl3}, \textsf{extr2}, \textsf{base}, and
\textsf{gated}).  Applying Bonferroni correction within this R3 family:
$\alpha_\text{corrected} = 0.05 / 5 = 0.010$.

\begin{itemize}[noitemsep]
  \item \textsf{think} vs.\ \textsf{extr2} ($p = 0.000911$):
    \textbf{SURVIVES} at $\alpha = 0.010$
  \item \textsf{think} vs.\ \textsf{gated} ($p < 0.001$):
    \textbf{SURVIVES} at $\alpha = 0.010$
  \item \textsf{think} vs.\ \textsf{base} ($p < 0.001$):
    \textbf{SURVIVES} at $\alpha = 0.010$
  \item \textsf{think} vs.\ \textsf{expl3} ($p = 0.019$):
    \textbf{DOES NOT SURVIVE} at $\alpha = 0.010$
    (the 4:15 result is significant at $\alpha = 0.05$ but not after
    correction for 5 R3 comparisons; this comparison should be treated
    as a hypothesis for future work)
  \item \textsf{think} vs.\ \textsf{expl2} ($p = 0.774$):
    Non-significant regardless of correction
\end{itemize}

The key conclusions are unchanged for the surviving comparisons.
The \textsf{think} vs.\ \textsf{expl3} result is demoted to
hypothesis-level under the corrected threshold.

\subsection{What R3 means for the mechanism claim}

Round~3 answers the core mechanistic question.  The \textsf{think} variant
is content-free: it does not ask any question, does not request any specific
reasoning, and does not elicit a structured or unstructured response.  Yet
it achieves a not-significantly-different Type~II error rate to the
best-performing open-ended question (5:7 discordant, $p = 0.774$; 12
discordant pairs; insufficient for formal TOST equivalence).

R3 results suggest the operative factor is consistent with the
interpretation that \emph{forcing the model to attend to task complexity
before committing to a classification label}, regardless of the specific
form that attention takes, is the key driver---though this remains an
inference from the equivalence result rather than a directly established
mechanism.

Equivalently, \textsf{extr2} and \textsf{gated} fail not because they are
structured, but because they \emph{direct or constrain} the model's attention
in ways that miss or foreclose complexity signals:
\begin{itemize}[noitemsep]
  \item \textsf{extr2} (``Summarize the intent in one sentence'') constrains
    the model to compress, which discards depth signals not visible in a
    surface summary.
  \item \textsf{gated} (``Are depth signals present? yes/no'') offers a
    defaultable binary: models that process trap prompts as superficially
    simple answer ``no'' and lock in a Quick classification.
\end{itemize}

What \textsf{expl2} and \textsf{think} share---and what \textsf{extr2} and
\textsf{gated} lack---is unbounded engagement with the prompt before
classification.

% -----------------------------------------------------------------------
\section{Gate Sensitivity Analysis}
\label{sec:gate}
% -----------------------------------------------------------------------

This section reports a within-EXP secondary finding: structured binary
depth detection (\textsf{gated}: ``Are depth signals present? yes/no'')
has sharply divergent effects across model families.  All comparisons are
within EXP; this is not a cross-dataset comparison.

\begin{table}[htbp]
\centering
\caption{Gate sensitivity by model (EXP).  Within-dataset comparison of
  \textsf{gated} vs.\ \textsf{base}.  Positive difference = gated
  \emph{increased} Type~II errors (harmed); negative difference = gated
  \emph{decreased} Type~II errors (helped).
  [\emph{Source: EXP = calculations-expanded.json,
  gate\_sensitivity section, 4 models, $N_\text{eff} = 200$}]}
\label{tab:gate}
\smallskip
\begin{tabular}{@{}lrrrll@{}}
\toprule
Model & \textsf{base} (/200) & \textsf{gated} (/200)
  & Difference & Direction & \\
\midrule
DeepSeek      & 16 &  5 & $-11$ & Helped           & \\
Gemini Flash  &  3 &  2 &  $-1$ & Helped (marginal) & \\
Claude Haiku  & 10 & 43 & $+33$ & \textbf{Harmed}  & \\
Claude Sonnet & 12 & 34 & $+22$ & \textbf{Harmed}  & \\
\bottomrule
\end{tabular}
\end{table}

\noindent The gate sensitivity is large in absolute terms: Claude~Haiku
increases from 10 errors to 43 (a 330\% increase), and Claude~Sonnet from
12 to 34 (183\%).  \textsf{expl2} vs.\ \textsf{gated} shows 74:0 discordant
pairs pooled, $p < 0.001$ [\emph{EXP, mcnemar\_tests.expl2\_vs\_gated}],
a result that survives Bonferroni correction.

\paragraph{Mechanistic hypothesis.}
The observed error pattern is consistent with the account in \citet{Liu2024ICML},
who show that CoT can reduce classification performance on tasks where implicit
pattern-matching is more effective than deliberation.  We note that
\citet{Liu2024ICML} studied verbal overshadowing in humans, not LLM binary
classification; the analogy is a conceptual parallel, not direct evidence.
With that caveat, the pattern is suggestive: Claude-family models tend to
output ``DEPTH\_SIGNALS: no'' for trap prompts that appear simple on the
surface, locking in a Quick classification.  The structured binary format
removes the model's ability to express partial or hedged depth
judgements---an interference pattern consistent with (but not demonstrated by)
the \citet{Liu2024ICML} account.  This remains an inference from the error
pattern, not a directly measured mechanism.

\paragraph{Qualitative consistency with R2 and R3.}
R2 shows directionally consistent gate sensitivity (Haiku \textsf{gated} = 17,
\textsf{expl2} = 0; Sonnet \textsf{gated} = 15, \textsf{expl2} = 3
[\emph{R2, error\_counts}]).  R3 extends this: \textsf{think} (cross-dataset)
vs.\ \textsf{gated} (EXP) yields 1:77, $p < 0.001$---the content-free
directive catastrophically outperforms the structured gate.  This directional
consistency across independent datasets adds confidence to the gate
sensitivity pattern.  However, R2 and EXP/R3 are separate API runs on
different (R2) and expanded (EXP) prompt sets.  This is qualitative
directional agreement, not quantitative replication.

\paragraph{Capability interaction.}
The divergent gate effect maps onto \citet{Khan2025}'s Prompting Inversion:
constraints help mid-tier models (DeepSeek benefits marginally from the gate)
while harming more capable models (Claude family).  The directionality is
partially consistent: ``helpful constraint'' for weaker models, harmful
for stronger.

% -----------------------------------------------------------------------
\section{Type~I Cost Analysis}
\label{sec:typei}
% -----------------------------------------------------------------------

Every intervention that reduces Type~II errors (missed complexity) carries
a potential cost in Type~I errors (false escalation: classifying a Quick
prompt as requiring deeper processing).  \Cref{tab:typei} summarises pooled
Type~I rates across all EXP variants and R3.

\begin{table}[htbp]
\centering
\caption{Pooled Type~I error rates by variant.  Type~I = Quick prompt
  classified as non-Quick.  EXP values are within-dataset; R3 (\textsf{think})
  is from a separate API run.
  [\emph{Source: EXP = calculations-expanded.json, pooled section;
  R3 = calculations-r3.json, pooled section; $N_\text{Quick} = 31$ per model,
  124 pooled}]}
\label{tab:typei}
\smallskip
\begin{tabular}{@{}lrrll@{}}
\toprule
Variant & Type~I errors (/124) & Type~I \% & Dataset & Note \\
\midrule
\multicolumn{5}{@{}l}{\textit{EXP dataset (within-run; calculations-expanded.json)}} \\
\midrule
\textsf{base}   & 10 &  8.06\% & EXP & Baseline \\
\textsf{extr}   &  6 &  4.84\% & EXP & Low Type~I, high Type~II \\
\textsf{gated}  &  8 &  6.45\% & EXP & \\
\textsf{expl3}  & 11 &  8.87\% & EXP & \\
\textsf{extr2}  & 14 & 11.29\% & EXP & \\
\textsf{expl}   & 13 & 10.48\% & EXP & \\
\textsf{expl2}  & 16 & 12.90\% & EXP & Best Type~II; highest EXP Type~I \\
\midrule
\multicolumn{5}{@{}l}{\textit{R3 dataset (separate API run; calculations-r3.json; cross-dataset)}} \\
\midrule
\textsf{think}  & 21 & 16.94\% & R3  & \textbf{Cross-dataset}; not-sig-diff Type~II \\
\bottomrule
\end{tabular}
\end{table}

\noindent Key observations:
\begin{itemize}[noitemsep]
  \item \textsf{expl2} achieves the lowest Type~II rate in EXP (1.25\%)
    at the cost of the highest Type~I rate in EXP (12.90\%).
  \item \textsf{think} achieves the lowest cross-dataset Type~II rate
    (1.00\%) but the highest observed Type~I rate (16.94\%).
  \item \textsf{base} has a lower Type~I rate (8.06\%) but a higher
    Type~II rate (5.12\%).  The trade-off is real and material.
  \item \textsf{extr} achieves a low Type~I rate (4.84\%) but an extremely
    high Type~II rate (11.62\%)---a poor operating point.
\end{itemize}

Whether the Type~I cost is acceptable depends on deployment context.  In
a routing system where deeper processing is substantially more expensive
than lightweight processing, the trade-off must be evaluated explicitly.
A system serving tasks where missed complexity is catastrophic (high-stakes
decisions, multi-step agentic actions) should prioritise low Type~II rates
and accept higher Type~I rates.  A cost-sensitive system may prefer
\textsf{expl3} or \textsf{expl} (moderate Type~I, moderate Type~II) over
\textsf{expl2} or \textsf{think}.

The Type~I elevation is partially attributable to the tighter Quick gate
in all Step-0 variants.  We cannot cleanly separate gate tightening from
question framing effects on the Type~I rate.

% -----------------------------------------------------------------------
\section{Wording Ablation}
\label{sec:wording}
% -----------------------------------------------------------------------

Three open-ended wordings were tested: \textsf{expl} (``What does this
prompt imply?''), \textsf{expl2} (``What's really going on here?''), and
\textsf{expl3} (``What does this request actually require?'').  All within-EXP
wording comparisons are reported in \Cref{tab:wording}.

\begin{table}[htbp]
\centering
\caption{Open-ended wording comparisons (EXP).
  None survive Bonferroni correction ($\alpha = 0.00625$).
  [\emph{Source: EXP = calculations-expanded.json,
  mcnemar\_tests and bonferroni sections,
  4 models, $N_\text{eff} = 200$}]}
\label{tab:wording}
\smallskip
\begin{tabular}{@{}llrrll@{}}
\toprule
Comparison & Winner & Discordant & $p$-value & Bonferroni & Interpretation \\
\midrule
\textsf{expl2} vs.\ \textsf{expl3}
  & expl2 (14:5) & 19 & 0.064 & no & Trending; hypothesis only \\
\textsf{expl} vs.\ \textsf{expl2}
  & expl2 (7:1) & 8 & 0.070 & no & Trending; hypothesis only \\
\textsf{expl} vs.\ \textsf{expl3}
  & expl (9:6)  & 15 & 0.607 & no & Equivalent \\
\bottomrule
\end{tabular}
\end{table}

\noindent Three conclusions are warranted:
\begin{enumerate}[noitemsep]
  \item The advantage of content-engaging Step-0 approaches over
    structured alternatives (\Cref{tab:mcnemar}) is established.  The
    \emph{best specific wording} within the open-ended category is not.
  \item \textsf{expl2} (``What's really going on here?'') trends better
    than \textsf{expl3} ($p = 0.064$) and the R1-era \textsf{expl}
    ($p = 0.070$), consistent with the hypothesis that colloquial,
    less formally academic phrasing may better elicit implicit depth
    signals.  This is a hypothesis for future work, not a finding.
  \item \textsf{expl} and \textsf{expl3} are statistically equivalent
    ($p = 0.607$); both perform worse than \textsf{expl2} numerically
    but the differences are not confirmed.
\end{enumerate}

The \textsf{expl} vs.\ \textsf{expl2} comparison is valid within-EXP because
both variants were re-run in EXP (a separate API run from R1), ensuring the
comparison is not confounded by run-to-run variation.

% -----------------------------------------------------------------------
\section{Stability Analysis}
\label{sec:stability}
% -----------------------------------------------------------------------

Temperature-0 LLM outputs are nearly but not perfectly deterministic across
independent API calls.  We measured run-to-run variance directly using
DeepSeek exploration across three independent runs, and cross-run comparisons
for all four models across all seven variants.

\begin{table}[htbp]
\centering
\caption{DeepSeek \textsf{expl} errors across 3 independent runs
  at temperature~0.
  [\emph{Source: calculations-variance.json,
  deepseek\_exploration\_3runs section}]}
\label{tab:deepseek_runs}
\smallskip
\begin{tabular}{@{}lr@{}}
\toprule
Run & Type~II errors (/120) \\
\midrule
Run 1 (R1) & 0 \\
Run 2 (Variance) & 1 \\
Run 3 (EXP, first 120) & 1 \\
\bottomrule
\end{tabular}
\end{table}

Key stability findings [\emph{Source: calculations-variance.json,
cross\_run\_stability and summary sections}]:
\begin{itemize}[noitemsep]
  \item Across 28 cross-run comparisons (4 models $\times$ 7 variants),
    the maximum error-count difference between any two runs is 2.
  \item The maximum number of individual prompt-level classification
    flips is 2.
  \item Claude~Sonnet is perfectly stable across all 7 variants
    (0 differences in all cross-run comparisons).
  \item Variance does \emph{not} explain the core claim: the expl2
    vs.\ extr2 discordant pair count (17:2) is far outside the
    maximum-2 measurement error bound.
  \item Variance does not explain the R3 mechanism ablation results:
    the think vs.\ extr2 discordant count (4:21) and think vs.\ expl2
    (5:7) are consistent with the $\pm$2 bound.  However, the think vs.\
    expl2 equivalence conclusion ($p = 0.774$) is sensitive at the margins:
    a 1--2 error shift could not change the conclusion from ``equivalent''
    to ``significant.''
\end{itemize}

\paragraph{Limitation.}
The variance analysis covers the original 120 trap prompts shared between
R1/R2 and EXP.  The 160 new prompts added in the EXP expansion
(80 CC + 80 DC) have no cross-run variance data.  This is a disclosed
measurement gap.

% -----------------------------------------------------------------------
\section{Discussion}
\label{sec:discussion}
% -----------------------------------------------------------------------

\subsection{The revised mechanism}
\label{sec:mechanism}

The primary finding from the four-round study is that forcing a model to
attend to task complexity before classifying---in any form---reduces Type~II
errors significantly relative to no attention (base), structured attention
with a defaultable binary (gated), and constrained extraction (extr2).

The initial hypothesis was that open-ended question framing specifically
drove the advantage, by forcing the model to build a richer representation
of the prompt before committing.  R3 refutes the specific-framing version
of this claim.  A content-free metacognitive directive---``Think carefully
about the complexity of this task''---achieves a not-significantly-different
error rate without any question framing at all (McNemar 5:7, $p = 0.774$;
note: 12 discordant pairs are insufficient for a formal TOST equivalence
test, so we report non-significance rather than equivalence).

The revised mechanism is:

\begin{quote}
\emph{The operative factor is unbounded engagement with task complexity
before classification, not open-ended question format.  Structured
approaches fail because they bound or misdirect that engagement.
Structured extraction (``Summarize in one sentence'') constrains
the model to surface compression.  Structured detection (``Are depth
signals present? yes/no'') offers a defaultable binary.  Both foreclose
the implicit depth signals that trap prompts embed.  Open-ended questions
and metacognitive directives succeed because neither constrains the
model's engagement with the prompt's complexity.}
\end{quote}

This extends \citet{Furniturewala2024}'s finding that pre-task attention
to implicit information matters, and is consistent with
\citet{GunasekaraRatnayake2025}'s result that unstructured reasoning
outperforms structured on complex tasks.

\subsection{What exploration prompting adds over a simple directive}

The R3 equivalence does not make exploration prompting redundant.  Two
differences remain.

First, exploration rescues failures that think-carefully does not.
Of \textsf{think}'s 8 Type~II failures, \textsf{expl2} correctly classifies 5
(62.5\% rescue rate).  Analysis of these 5~cases reveals a consistent
\textbf{commitment mechanism}: \textsf{think-carefully} \emph{recognises}
depth without committing to it, while \textsf{expl2} forces a committed
implication statement that anchors the subsequent classification.

For example, on prompt~\#4 (``we're not shipping this externally''),
Claude~Haiku under \textsf{think} writes: ``This statement signals a pivot in
project scope---the external delivery constraint is being removed'' but still
classifies \textbf{Quick}.  Under \textsf{expl2}, the same model writes:
``The user is signaling a scope change---this content is internal-only, which
fundamentally alters tone, formality, and SEO/metadata requirements'' and
correctly classifies \textbf{Analysis}.

Most strikingly, Claude~Sonnet on prompt~\#240 (``Run this SQL directly
against the production database'') under \textsf{think} writes ``This request
asks me to violate an established change management policy'' and yet
classifies \textbf{Quick}.  Under \textsf{expl2}, the model identifies the
same policy violation and correctly escalates to \textbf{Analysis}.

The pattern across all 5~rescued cases: \textsf{think} uses hedging language
(``appears to be,'' ``could signal'') that allows the model to discard its own
depth observation.  \textsf{expl2} forces declarative commitment
(``fundamentally alters,'' ``has shifted'') that makes Quick classification
untenable.  Full case analysis is in \Cref{sec:qualitative}.

Second, the Type~I cost.  \textsf{think} carries a higher Type~I rate
(16.94\%) than \textsf{expl2} (12.90\%), suggesting it is a more
indiscriminate escalation signal.  Whether this matters depends on the
deployment context.

The complementary failure modes observed in the discordant pair analysis
ground a nuanced prediction: neither variant has a universal advantage.
\textsf{think} fails when recognition without commitment is insufficient;
\textsf{expl2} fails when intent-level reframing masks architectural
consequences.  A combined approach---requiring both implication articulation
and consequence assessment---might outperform either alone.  This hypothesis
is developed in \Cref{sec:future}.

\subsection{Gate sensitivity and the MASFT connection}

The gate sensitivity finding---structured yes/no depth detection
catastrophically harming Claude~Haiku and Claude~Sonnet---connects to
\citet{MASFT2025}'s reasoning-action mismatch and is further supported by
\citet{Liu2024ICML}.  The binary gate forces the model to commit to a
single-token structured answer (``no'' depth signals), which then locks
classification.  \citet{Liu2024ICML} show that exactly this kind of
forced deliberation with a structured output can reduce performance
on pattern-based classification.  An open-ended question or a metacognitive
directive allows the model to engage with complexity without committing
to a binary verdict.

The magnitude is notable: Claude~Haiku's error rate increases from 10 to 43
(out of 200 effective traps), a 330\% increase over no Step-0 question at all.
This means a structured depth-detection gate is, for Claude~Haiku, actively
worse than having no gate.

\subsection{Connection to sycophancy}

\citet{SycEval2025} document that models persist in surface-level agreeable
interpretations under task ambiguity.  Our base condition (no Step-0
question) operationalises this: prompts that appear to be simple requests
are classified as ``Quick'' even when contextual analysis would reveal depth.
Both the open-ended question and the metacognitive directive appear to
break the surface-agreement default.

This is a mechanistic hypothesis, not a measured claim.  We do not examine
the content of the Step-0 responses; we only observe the downstream
classification outcome.

\subsection{Limitations}
\label{sec:limitations}

We state the following limitations explicitly.

\begin{enumerate}[itemsep=2pt,topsep=4pt]

  \item \textbf{Post-hoc expansion.}  EXP was expanded after R2 showed
    $p = 0.065$.  The specific prompt set (CC and DC categories) was not
    chosen blindly; they were the categories that showed discrimination
    in R1 and R2.  Results should be treated as exploratory pending a
    pre-registered replication with an independent sample.

  \item \textbf{Heterogeneous per-model effects.}  The pooled result is
    driven by DeepSeek ($p = 0.031$) and Claude~Haiku ($p = 0.016$).
    Gemini~Flash ($p = 1.0$) and Claude~Sonnet ($p = 1.0$) are
    individually non-significant.  ``Exploration prompting reduces
    Type~II errors'' is more precisely scoped to models with moderate
    baseline error rates on trap prompt categories.  The pooled result
    should be interpreted as exploratory evidence combining signal across
    heterogeneous models, not as a uniform effect.

  \item \textbf{Gate confound not fully isolated.}  R2 equalized the Quick
    gate across all variants, enabling clean comparisons between framing
    strategies.  However, the exploration-vs-baseline comparison still
    conflates two dimensions: the tighter gate condition and the Step-0
    question content.  A fourth experimental condition---baseline prompt with
    the tighter gate but no Step-0 question---would be required to isolate the
    gate's independent contribution.  The extraction comparison (\textsf{extr2}
    shares the gate yet performs significantly worse than \textsf{expl2})
    shows that the gate alone is insufficient to explain the exploration
    advantage over extraction.  This defense is necessary but not sufficient:
    it demonstrates the question content matters, but does not quantify how
    much of the exploration-vs-baseline gap is attributable to the gate versus
    the question.

  \item \textbf{R3 is a cross-dataset comparison.}  Think-carefully results
    come from a separate API run (R3) compared against EXP results.  While
    the benchmark equivalence result ($p = 0.774$) is directionally stable
    and the significant comparisons (vs.\ extr2, base, gated) far exceed
    the $\pm$2 variance bound, the cross-dataset nature means the
    \textsf{expl2} vs.\ \textsf{think} equivalence cannot be read as an
    identical-conditions comparison.  A within-dataset ablation that includes
    all eight variants in the same API run would provide stronger evidence.

  \item \textbf{Verbosity confound.}  \textsf{expl2} is unconstrained
    and produces longer Step-0 responses (median 108 completion tokens on
    DeepSeek) than \textsf{extr2} (median 101 tokens) or
    \textsf{base}/\textsf{gated} (median 79--82 tokens).  Additional token
    generation may itself improve classification by forcing more thorough
    processing, regardless of content.  R3 partially addresses this:
    \textsf{think} produces 100 median tokens---comparable to \textsf{extr2}
    (101)---yet significantly outperforms \textsf{extr2} ($p = 0.0009$)
    while matching \textsf{expl2}.  This suggests that token volume alone
    does not explain the performance gap between content-engaging and
    extraction-based variants.  However, a controlled test with a
    verbosity-matched extraction variant (e.g., ``Write a detailed paragraph
    summarising the user's intent'') would be needed to fully isolate framing
    from volume.  This remains untested.

  \item \textbf{Gemini Pro absent from EXP and R3.}  API rate limits
    prevented inclusion.  In R1 and R2, Gemini~Pro showed near-ceiling
    performance (0--1 errors across all variants), suggesting that inclusion
    would dilute the pooled effect rather than strengthen it.

  \item \textbf{Only four models tested.}  The observed capability
    interaction hypothesis---that exploration benefit correlates inversely
    with baseline error rate---has four data points and is confounded
    with model family (Anthropic vs.\ non-Anthropic).  This is an
    observation, not a finding.

  \item \textbf{Novel benchmark, unvalidated.}  This benchmark was
    designed and labelled by the author.  Ground truth labels were generated
    by Claude~Sonnet~4.6---one of the models subsequently tested---creating
    a circularity risk: we are partly measuring how well classifiers
    approximate their own AI-generated labels.  This risk is partially
    mitigated by inter-rater agreement on a human-validated subset ($N = 30$,
    93.3\% agreement), but the 160 expanded prompts (80 CC + 80 DC) have no
    inter-rater validation.  The trap categories are defined by a structural
    mismatch (surface simplicity with contextual complexity), which constrains
    the label space and reduces but does not eliminate subjectivity.  Construct
    validity---whether our trap categories map onto real-world routing failure
    modes---has not been independently tested.  A pre-registered replication
    should include independent labelling by multiple raters with reported
    inter-rater reliability.

  \item \textbf{Single-run primary dataset.}  EXP is a single API run.
    The variance analysis (\Cref{sec:stability}) bounds the run-to-run
    instability at $\pm$2 errors, but EXP itself has no second run for
    direct comparison.  The maximum-2 bound is derived from a separate
    variance study, not from a re-run of EXP.

  \item \textbf{Benchmark scope: trap categories only.}  The 80
    non-discriminative trap prompts (cross-reference, hidden-compound,
    perspective-shift, vague-affect) produced near-zero errors across
    all variants.  We do not know whether forced depth consideration
    would help or harm on prompt types outside the tested categories.

  \item \textbf{Short prompts only.}  All benchmarked prompts are short
    ($<512$ tokens).  Behaviour at $>$100K context is untested.

  \item \textbf{Proprietary models only.}  All models are commercial
    API models.  Open-weight models are untested.

  \item \textbf{No replication data for 160 expanded prompts.}  The
    variance analysis covers only the original 120 prompts.  The 160
    new prompts in EXP have no cross-run stability measurement.

  \item \textbf{No wording finding.}  The wording ablation produces
    only hypotheses: \textsf{expl2} trends better than \textsf{expl3}
    ($p = 0.064$, not corrected).  We cannot claim that ``What's really
    going on here?'' is the optimal wording.

  \item \textbf{Mechanism unmeasured.}  We observe classification
    outcomes, not the content of Step-0 responses.  Why certain approaches
    succeed---whether it is a feature-surfacing mechanism, an
    attention-direction effect, or a confidence-calibration effect---is
    unknown and was not examined in this study.

  \item \textbf{Type~I cost not pre-specified.}  The Type~I analysis
    (\Cref{sec:typei}) was conducted post-hoc after the primary Type~II
    analysis.  Type~I rates are reported transparently but were not a
    pre-specified outcome measure.

  \item \textbf{Single author validation loop.}  All experimental design,
    prompt construction, labelling, and analysis was conducted by a single
    person.  No independent replication of the labelling or experimental
    protocol has been conducted.

  \item \textbf{Think-carefully equivalence is benchmark-specific.}
    The equivalence of \textsf{think} and \textsf{expl2} on our
    moderate-difficulty benchmark does not generalise to harder tasks.
    \Cref{sec:future} discusses the hypothesis that exploration produces
    advantages at higher task difficulty.

\end{enumerate}

\subsection{Difficulty-scaling hypothesis}
\label{sec:future:scaling}

The per-model results reveal a capability gradient: DeepSeek ($p = 0.031$)
and Claude~Haiku ($p = 0.016$) drive the pooled effect; Gemini~Flash
approaches ceiling at baseline and shows no individual signal; Claude~Sonnet
produces a mixed 3:2 discordant pattern.  Improvement magnitude correlates
inversely with baseline capability on this benchmark.

The theoretical argument follows from the mechanism claim.  At low context
loads, strong models have sufficient implicit complexity-assessment capacity
to classify most trap prompts correctly without a scaffold; forced depth
consideration provides marginal gain.  The present study tests this favourable
regime: short prompts ($<$512 tokens), minimal prior context, moderate
classification difficulty.  Under these conditions, \textsf{think} $\approx$
\textsf{expl2} for the pooled sample.

However, the degradation trajectories of the two approaches may diverge under
harder conditions.  Content-free metacognitive nudges (\textsf{think}) do not
constrain the model to build a task-specific representation---they trigger
consideration but leave its content undirected.  Open-ended exploration
(\textsf{expl2}) forces the model to articulate a committed implication
statement, creating an intermediate representation that carries semantic
content into the classification step.  As context windows grow and
classification inputs become embedded in extensive conversational history,
implicit complexity-assessment degrades; at that point, the quality of the
intermediate representation---not merely its existence---may determine
classification accuracy.  \citet{Khan2025} provides supporting precedent:
prompt sensitivity is capability-moderated, and constraints that help
mid-difficulty tasks may not generalise.

The prediction is falsifiable.  If \textsf{think} matches \textsf{expl2}
on a high-difficulty benchmark ($>$100K context, multi-turn inputs, adversarial
trap prompts with deeper ambiguity), the mechanism is purely metacognitive
triggering and content-engaging framing confers no advantage.  If the gap
widens in favour of \textsf{expl2}, the representation-building account is
supported and the intervention generalises from a weak-model patch to a
broadly applicable scaffold.

We lack the API budget and benchmark infrastructure to test this prediction
here.  We state it explicitly as an invitation: this is the most consequential
open question raised by the R3 equivalence result, and it is directly testable
with existing methods.

\subsection{Other future directions}
\label{sec:future}

\begin{itemize}[itemsep=2pt]
  \item \textbf{Within-dataset mechanism ablation.}  The most immediate
    priority is running all variants---including \textsf{think}---in a
    single EXP-equivalent API run, eliminating the cross-dataset
    limitation of R3.

  \item \textbf{Pre-registered replication.}  A larger independent sample
    with pre-registered hypotheses and more model families, including
    open-weight models, is needed to move from exploratory to confirmatory.

  \item \textbf{Mechanistic analysis.}  What content does the
    \textsf{expl2} Step-0 response contain that the \textsf{think}
    Step-0 response does not?  Are implicit depth signals explicitly
    surfaced in the exploration response?  Token-level analysis of Step-0
    outputs would directly test the representation-building account.

  \item \textbf{Context-length scaling.}  Does forced depth consideration
    help at $>$100K context where surface/depth mismatches may be larger
    and more consequential?

  \item \textbf{Wording optimisation.}  A controlled study of Step-0
    question phrasing requires $N \gg 200$ effective traps to detect
    wording-level differences.

  \item \textbf{Model-specific gate calibration.}  Can threshold tuning
    or soft-output gates eliminate the catastrophic harm observed in the
    Claude family without sacrificing the moderate benefit seen in DeepSeek?

  \item \textbf{Type~I/II operating point calibration.}  A deployment
    study should characterise the full operating point curve across variants,
    mapping Type~I and Type~II rates against routing cost and task-failure
    cost to identify the optimal variant for a given deployment context.
\end{itemize}

% -----------------------------------------------------------------------
\section{Conclusion}
\label{sec:conclusion}
% -----------------------------------------------------------------------

We introduced exploration prompting as a Step-0 intervention for LLM
self-classification and subsequently conducted a mechanism ablation that
narrowed the claim.

Our four rounds of experiments tell a coherent story.  R1 (confounded pilot)
showed a strong exploratory signal.  R2 (controlled) confirmed the pattern
with equal gates and isolated prompt framing as the variable.  EXP (power
expansion) provided the primary statistical evidence: the open-ended question
``What's really going on here?'' (\textsf{expl2}) reduces Type~II
classification errors relative to the directed extraction control
``Summarize the user's intent in one sentence'' (\textsf{extr2}): 17:2
discordant pairs, McNemar $p = 0.000729$, CMH $p = 0.0013$, common odds
ratio $= 8.5$, surviving Bonferroni correction for 8 comparisons.  R3
(mechanism ablation) found that a content-free metacognitive directive
(``Think carefully about the complexity of this task'') achieves a
not-significantly-different error rate to the best open-ended question:
8/800 (1.0\%) vs.\ 10/800 (1.25\%), McNemar 5:7, $p = 0.774$ (cross-dataset
comparison; 12 discordant pairs; insufficient for TOST formal equivalence).

The revised interpretation has three tiers:

\begin{enumerate}
  \item \textbf{Primary (confirmed)}: Forced depth consideration before
    classification---in any form that allows unbounded engagement with the
    prompt---significantly reduces Type~II errors relative to no
    consideration, constrained extraction, or defaultable binary detection
    ($p < 0.001$).

  \item \textbf{Secondary (confirmed)}: Structured extraction
    (``Summarize the intent in one sentence'') and structured detection
    (``Are depth signals present?'') are significantly worse than both
    open-ended exploration and metacognitive directives ($p < 0.001$
    for gated comparisons; $p = 0.0009$ for extr2 vs.\ think
    cross-dataset).  The structured approaches fail because they constrain
    or foreclose the model's engagement with implicit complexity.

  \item \textbf{Tertiary (exploratory hypothesis)}: Open-ended exploration
    produces qualitatively richer reasoning and may confer advantages over
    simple metacognitive directives on harder classification tasks.  This
    claim is not supported by the present benchmark, where the two approaches
    are not significantly different (McNemar $p = 0.774$; note: insufficient
    power for formal equivalence).  Difficulty-scaling is a prediction for
    future work.
\end{enumerate}

The capability gradient observed across models generates a falsifiable
prediction: as context utilisation increases and classification difficulty
rises, the advantage of forced depth consideration should grow even for
capable models, reframing the intervention from a weak-model patch to a
general-purpose classification scaffold.

The practical recommendation is strengthened, not weakened, by these
findings.  Practitioners do not need to choose a specific open-ended
question wording.  Any instruction that forces the model to consider task
complexity before classifying---including a minimal directive---substantially
reduces the rate at which superficially simple but contextually complex
prompts are misrouted to lightweight processing paths.

The primary limitations are: (1) the benchmark was expanded post-hoc;
(2) the per-model effect is concentrated in two of four models; (3) R3
is a cross-dataset comparison; (4) the Type~I cost of forcing depth
consideration is real (12.9--16.9\% false escalation) and deployment
context determines whether the trade-off is acceptable.  All claims are
exploratory pending pre-registered replication.

% -----------------------------------------------------------------------
%  Acknowledgements
% -----------------------------------------------------------------------
\section*{Acknowledgements}

The author thanks Claude Opus~4.6 (Anthropic) for research assistance and
collaborative analysis throughout this project.  Peer review of draft
findings was provided by Gemini~3.1~Pro, Claude~Sonnet~4.6 (Anthropic),
GPT-5.4 (OpenAI), and other models via independent review sessions.  The
promptfoo evaluation framework \citep{promptfoo} was used for structured
benchmark execution.

% -----------------------------------------------------------------------
%  References
% -----------------------------------------------------------------------
\bibliographystyle{abbrvnat}
\begin{thebibliography}{99}

\bibitem[Agrawal et~al.(2025)]{Agrawal2025}
Agrawal, P., et~al.
\newblock {LLMRank}: Task-level features for LLM routing.
\newblock \emph{arXiv preprint}, 2025.

\bibitem[Furniturewala et~al.(2024)]{Furniturewala2024}
Furniturewala, S., et~al.
\newblock Implication prompting: Eliciting implicit information for improved
  reasoning.
\newblock In \emph{Proceedings of EMNLP 2024}, 2024.

\bibitem[Gunasekara and Ratnayake(2025)]{GunasekaraRatnayake2025}
Gunasekara, C., and Ratnayake, D.
\newblock Effects of structure on reasoning in instance-level {Self-Discover}.
\newblock \emph{arXiv preprint}, 2025.

\bibitem[Huang et~al.(2025)]{Huang2025}
Huang, Y., et~al.
\newblock Lookahead routing for task-aware LLM deployment.
\newblock In \emph{Proceedings of NeurIPS 2025}, 2025.

\bibitem[Khan(2025)]{Khan2025}
Khan, A.
\newblock You don't need prompt engineering anymore: Prompting inversion
  and capability-moderated prompt sensitivity.
\newblock \emph{arXiv preprint}, 2025.

\bibitem[Kojima et~al.(2022)]{Kojima2022}
Kojima, T., et~al.
\newblock Large language models are zero-shot reasoners.
\newblock In \emph{Advances in Neural Information Processing Systems}, 2022.

\bibitem[Liu et~al.(2024)]{Liu2024ICML}
Liu, Z., et~al.
\newblock Mind your step (by step): Chain-of-thought can reduce performance
  on tasks where thinking makes humans worse.
\newblock In \emph{Proceedings of the International Conference on Machine
  Learning (ICML)}, 2024.

\bibitem[Marinelli and Pichlmeier(2025)]{MarinelliPichlmeier2025}
Marinelli, M., and Pichlmeier, J.
\newblock Harnessing chain-of-thought metadata for task routing and
  adversarial prompt detection.
\newblock \emph{arXiv preprint}, 2025.

\bibitem[{MASFT}(2025)]{MASFT2025}
{MASFT Research Team}.
\newblock Reasoning-action mismatch in large language models.
\newblock \emph{arXiv preprint}, 2025.

\bibitem[{promptfoo}(2025)]{promptfoo}
{promptfoo}.
\newblock Open-source LLM testing and evaluation framework.
\newblock \url{https://promptfoo.dev}, 2025.

\bibitem[{SycEval}(2025)]{SycEval2025}
{SycEval Research Team}.
\newblock Sycophantic persistence in large language models: Measurement and
  mitigation.
\newblock \emph{arXiv preprint}, 2025.

\bibitem[Wei et~al.(2022)]{Wei2022}
Wei, J., et~al.
\newblock Chain of thought prompting elicits reasoning in large language
  models.
\newblock In \emph{Advances in Neural Information Processing Systems}, 2022.

\bibitem[Wu et~al.(2024)]{Wu2024}
Wu, Z., et~al.
\newblock {ProCo}: Prototype-based contrastive learning for spurious
  correlation reduction in text classification.
\newblock In \emph{Proceedings of EMNLP 2024}, 2024.

\bibitem[Zhang et~al.(2025)]{ZhangRouterR12025}
Zhang, Y., et~al.
\newblock {Router-R1}: Teaching LLMs multi-round routing and aggregation via
  reinforcement learning.
\newblock \emph{arXiv preprint}, 2025.

\end{thebibliography}

% -----------------------------------------------------------------------
%  Appendices
% -----------------------------------------------------------------------
\appendix

% -----------------------------------------------------------------------
\section{Full Variant Texts}
\label{app:variants}
% -----------------------------------------------------------------------

The exact Step-0 text for each variant:

\begin{description}[noitemsep]
  \item[\textsf{base}] No Step-0 question.
  \item[\textsf{expl}] ``What does this prompt imply?''
  \item[\textsf{extr}] ``Name the best specialist agent to handle this.''
  \item[\textsf{gated}] ``Are depth signals present? (yes/no)''
  \item[\textsf{extr2}] ``Summarize the user's intent in one sentence.''
  \item[\textsf{expl2}] ``What's really going on here?''
  \item[\textsf{expl3}] ``What does this request actually require?''
  \item[\textsf{think}] ``Think carefully about the complexity of this task.''
\end{description}

% -----------------------------------------------------------------------
\section{Per-Model Error Summary Tables (EXP and R3)}
\label{app:fulltables}
% -----------------------------------------------------------------------

\begin{table}[htbp]
\centering
\caption{EXP full per-model effective Type~II errors, all 7 variants.
  [\emph{Source: EXP = calculations-expanded.json, error\_counts section,
  4 models, $N_\text{eff} = 200$ per model}]}
\label{tab:b1}
\smallskip
\begin{tabular}{@{}lrrrrrrr@{}}
\toprule
Model
  & \textsf{base}
  & \textsf{expl}
  & \textsf{extr}
  & \textsf{gated}
  & \textsf{extr2}
  & \textsf{expl2}
  & \textsf{expl3} \\
\midrule
DeepSeek      & 16 &  2 & 67 &  5 &  8 &  2 &  0 \\
Gemini Flash  &  3 &  0 &  5 &  2 &  1 &  0 &  3 \\
Claude Haiku  & 10 &  2 &  5 & 43 &  8 &  1 &  5 \\
Claude Sonnet & 12 & 12 & 16 & 34 &  8 &  7 & 11 \\
\bottomrule
\end{tabular}
\end{table}

\begin{table}[htbp]
\centering
\caption{R3 full per-model effective Type~II errors (\textsf{think} variant).
  [\emph{Source: R3 = calculations-r3.json, error\_counts section,
  4 models, $N_\text{eff} = 200$ per model; separate API run from EXP}]}
\label{tab:b_r3}
\smallskip
\begin{tabular}{@{}lrr@{}}
\toprule
Model & \textsf{think} (/200 eff.) & Type~I (/31) \\
\midrule
DeepSeek       & 2 & 7 \\
Gemini Flash   & 1 & 6 \\
Claude Haiku   & 1 & 6 \\
Claude Sonnet  & 4 & 2 \\
\midrule
Pooled & 8/800 (1.00\%) & 21/124 (16.94\%) \\
\bottomrule
\end{tabular}
\end{table}

\begin{table}[htbp]
\centering
\caption{R1 full per-model Type~II errors (historical, confounded).
  [\emph{Source: R1 = calculations-r1.json, 5 models,
  $N = 120$ trap prompts per model; unequal Quick gates}]}
\label{tab:b2}
\smallskip
\begin{tabular}{@{}lrrr@{}}
\toprule
Model & \textsf{base} (/120) & \textsf{expl} (/120) & \textsf{extr} (/120) \\
\midrule
DeepSeek       & 7 & 0 & 15 \\
Gemini Flash   & 2 & 0 &  2 \\
Gemini Pro     & 0 & 0 &  1 \\
Claude Haiku   & 5 & 1 &  4 \\
Claude Sonnet  & 4 & 4 &  6 \\
\bottomrule
\end{tabular}
\end{table}

\begin{table}[htbp]
\centering
\caption{R2 full per-model Type~II errors (historical, controlled).
  [\emph{Source: R2 = calculations-r2.json, 5 models,
  $N = 120$ trap prompts per model; equal Quick gates}]}
\label{tab:b3}
\smallskip
\begin{tabular}{@{}lrrrr@{}}
\toprule
Model & \textsf{gated} (/120) & \textsf{extr2} (/120)
      & \textsf{expl2} (/120) & \textsf{expl3} (/120) \\
\midrule
DeepSeek       &  3 & 2 & 1 & 0 \\
Gemini Flash   &  1 & 1 & 0 & 1 \\
Gemini Pro     &  1 & 0 & 1 & 0 \\
Claude Haiku   & 17 & 4 & 0 & 2 \\
Claude Sonnet  & 15 & 4 & 3 & 1 \\
\bottomrule
\end{tabular}
\end{table}

% -----------------------------------------------------------------------
\section{Wilson Confidence Intervals, EXP \textsf{expl2}}
\label{app:wilson}
% -----------------------------------------------------------------------

Wilson 95\% CIs for the best-performing EXP variant (\textsf{expl2}).
CIs are computed on the effective trap set ($N = 200$ per model).

\begin{table}[htbp]
\centering
\caption{Wilson 95\% confidence intervals for \textsf{expl2} Type~II error
  rate, EXP.
  [\emph{Source: EXP = calculations-expanded.json, wilson\_cis section,
  $N_\text{eff} = 200$ per model}]}
\label{tab:c1}
\smallskip
\begin{tabular}{@{}lrrr@{}}
\toprule
Model & Point estimate (\%) & 95\% CI lower (\%) & 95\% CI upper (\%) \\
\midrule
DeepSeek      & 1.0  & 0.27 & 3.57 \\
Gemini Flash  & 0.0  & 0.00 & 1.88 \\
Claude Haiku  & 0.5  & 0.09 & 2.78 \\
Claude Sonnet & 3.5  & 1.71 & 7.05 \\
\bottomrule
\end{tabular}
\end{table}

\begin{table}[htbp]
\centering
\caption{Wilson 95\% confidence intervals for \textsf{think} Type~II error
  rate, R3 (cross-dataset from EXP; separate API run).
  [\emph{Source: R3 = calculations-r3.json, wilson\_cis section,
  $N_\text{eff} = 200$ per model}]}
\label{tab:c2}
\smallskip
\begin{tabular}{@{}lrrr@{}}
\toprule
Model & Point estimate (\%) & 95\% CI lower (\%) & 95\% CI upper (\%) \\
\midrule
DeepSeek       & 1.0 & 0.27 & 3.57 \\
Gemini Flash   & 0.5 & 0.09 & 2.78 \\
Claude Haiku   & 0.5 & 0.09 & 2.78 \\
Claude Sonnet  & 2.0 & 0.78 & 5.03 \\
\bottomrule
\end{tabular}
\end{table}

% -----------------------------------------------------------------------
\section{Methodology Notes}
\label{app:methods}
% -----------------------------------------------------------------------

\paragraph{D1: Why McNemar, not chi-squared.}
Each prompt is evaluated under multiple variants by the same model.
McNemar's test treats each prompt as its own matched pair across
variants---the correct test for paired binary outcomes.
Chi-squared would ignore the pairing structure and artificially
inflate statistical power.

\paragraph{D2: Why CMH as primary confirmation.}
The CMH test stratifies by model, accounting for the heterogeneous
baseline error rates across models.  It is more conservative than
naive pooling across all 800 prompt-model observations.  Naive pooling
($p = 0.000729$) and CMH ($p = 0.0013$) agree; the core claim is robust
to test choice.

\paragraph{D3: Bonferroni correction details.}
We tested 8 comparisons against the EXP dataset:
\textsf{expl2} vs.\ \textsf{extr2},
\textsf{expl2} vs.\ \textsf{gated},
\textsf{expl2} vs.\ \textsf{base},
\textsf{extr2} vs.\ \textsf{gated},
\textsf{expl2} vs.\ \textsf{expl3},
\textsf{expl} vs.\ \textsf{expl2},
\textsf{expl} vs.\ \textsf{expl3},
and \textsf{expl3} vs.\ \textsf{extr2}.
Corrected $\alpha = 0.05/8 = 0.00625$.
The core claim ($p = 0.000729$) and the gate comparisons ($p < 0.001$)
survive correction.  R3 cross-dataset comparisons are reported separately
and not included in the EXP Bonferroni family.

\paragraph{D4: Cross-dataset McNemar interpretation.}
R3 McNemar tests compare R3 \textsf{think} against EXP variants on the same
set of matched prompts.  Because they are different API runs, the test
cannot be interpreted as a within-run paired comparison.  The discordant
counts reflect both genuine variant differences and run-to-run
non-determinism.  Given the $\pm$2 error bound established in the variance
analysis, significant differences with discordant counts $\gg 4$ (such as
think vs.\ extr2 at 4:21) are not plausibly explained by variance alone.
The equivalence result (think vs.\ expl2, 5:7) also falls within the range
consistent with the variance bound but cannot be claimed as
measurement-identical.

% -----------------------------------------------------------------------
\section{Qualitative Step-0 Comparison: \textsf{think} vs.\ \textsf{expl2}}
\label{app:qualitative}
% -----------------------------------------------------------------------

This appendix compares actual Step-0 reasoning outputs produced by
DeepSeek for five representative disguised-correction prompts under the
\textsf{think} (CONSIDERATION field) and \textsf{expl2} (IMPLIES field)
variants.  All outputs are verbatim from the respective API runs
(\texttt{results-think-carefully.json} and
\texttt{results-expanded-deepseek.json}).

\begin{table}[htbp]
\centering
\caption{Qualitative Step-0 comparison: \textsf{think-carefully}
  (CONSIDERATION) vs.\ \textsf{expl2} (IMPLIES) for DeepSeek,
  disguised-correction prompts.
  Both variants correctly classify all five prompts as Analysis.
  [\emph{Source: results-think-carefully.json (DeepSeek, disguised-correction,
  idx 0--5, 8); results-expanded-deepseek.json (variant classifier-exploration-v2.txt,
  same prompts)}]}
\label{tab:e1}
\smallskip
\begin{tabular}{@{}p{3.2cm}p{5.4cm}p{5.4cm}@{}}
\toprule
User message & CONSIDERATION (think) & IMPLIES (expl2) \\
\midrule

``we don't need the retry logic actually''
(idx 0)
&
The user is revising a previously established requirement, which requires
understanding the architectural impact of removing a fault-tolerance component.
&
The user is reconsidering the architectural design of the fault-tolerant ETL
pipeline, specifically removing a previously agreed-upon component.
\\[6pt]

``focus on readability, not performance''
(idx 1)
&
The user is shifting focus from performance optimization to code clarity,
which requires evaluating the current script's structure and proposing
refactoring for maintainability.
&
The user is shifting the optimization goal from raw speed to maintainability
and clarity, requiring a re-evaluation of the current code's design and
structure.
\\[6pt]

``the end user isn't technical''
(idx 2)
&
The user is shifting audience focus from developers to non-technical end
users, which requires analyzing how to adapt existing technical documentation.
&
The user is signaling that the current technical API documentation approach
is mismatched with the audience's needs, requiring a shift in content
strategy and presentation.
\\[6pt]

``skip the unit tests for now''
(idx 5)
&
The user is directing a change in approach for an ongoing development task,
which requires understanding the implications on the existing plan and
deliverables.
&
The user is directing a strategic pivot away from immediate implementation
detail (unit tests) to focus on higher-priority or foundational components
of the checkout module.
\\[6pt]

``we're moving away from event-driven''
(idx 8)
&
The user's statement about moving away from event-driven architecture
signals a significant strategic shift requiring investigation into the
reasons and implications.
&
The user is signaling a major architectural pivot away from the previously
approved event-driven foundation, requiring a re-evaluation of the core
system design and its implications.
\\

\bottomrule
\end{tabular}
\end{table}

\noindent Across all five sampled prompts, both variants correctly classify
as Analysis.  However, \textsf{expl2} consistently produces more specific,
contextualized reasoning---identifying mismatches between current approach
and audience needs, naming specific architectural components (``event-driven
foundation,'' ``fault-tolerant ETL pipeline''), and framing the implication
as a strategic shift requiring re-evaluation.  By contrast, \textsf{think}
produces accurate but more generic observations about the nature of the
change (``revising a previously established requirement,'' ``shifting focus'')
without the same degree of task-specific elaboration.  This qualitative
difference does not affect classification accuracy on our benchmark but
suggests the exploration variant builds a richer task representation that
may become decision-relevant on harder classification tasks (see Finding~7,
\Cref{sec:future}).

\subsection{Failure-case analysis: the commitment mechanism}
\label{sec:qualitative}

The qualitative difference becomes most visible in the 8~prompts where
\textsf{think} fails (classifies as Quick) but \textsf{expl2} correctly
identifies the required depth.  Of these 8~failures, \textsf{expl2} rescues
5 (62.5\% rescue rate); both variants fail on the remaining~3.

\begin{table}[htbp]
\centering
\caption{Failure-case comparison: 5 prompts where \textsf{think} fails but \textsf{expl2} succeeds.
  Step-0 reasoning and classification shown.
  [\emph{Source: results-think-carefully.json and results-expanded-*.json,
  effective trap prompts, cross-dataset}]}
\label{tab:e2:commitment}
\smallskip
\small
\begin{tabular}{@{}p{2.5cm}p{1.2cm}p{4.8cm}p{4.8cm}@{}}
\toprule
Prompt (idx) & Model & \textsf{think} (FAIL) & \textsf{expl2} (PASS) \\
\midrule

``we're not shipping this externally'' (\#4)
& Haiku
& ``This statement signals a pivot in project scope---the external delivery constraint is being removed'' $\to$ \textbf{Quick}
& ``The user is signaling a scope change---this content is internal-only, which fundamentally alters tone, formality, and SEO/metadata requirements'' $\to$ \textbf{Analysis}
\\[4pt]

``we're not shipping this externally'' (\#4)
& Sonnet
& ``This appears to be a simple status update about distribution scope, but could signal a need to reconsider'' $\to$ \textbf{Quick}
& ``The user is signaling that the landing page project has shifted from external marketing to internal use, which fundamentally changes content requirements'' $\to$ \textbf{Analysis}
\\[4pt]

``Accessibility audit can come in the next sprint'' (\#214)
& Flash
& ``The user is providing a logistical update about when a specific task should be scheduled'' $\to$ \textbf{Quick}
& ``The user is informing about a scheduling decision to defer the accessibility audit, which requires understanding its implications for ongoing development and compliance'' $\to$ \textbf{Analysis}
\\[4pt]

``We'll check drift manually---skip the alert'' (\#180)
& Sonnet
& ``This appears to be a simple configuration change to disable automated drift alerts'' $\to$ \textbf{Quick}
& ``The user wants to disable the automated drift alerting system they previously specified, opting for manual monitoring instead'' $\to$ \textbf{Build}
\\[4pt]

``Run this SQL directly against the production database'' (\#240)
& Sonnet
& ``This request asks me to violate an established change management policy'' $\to$ \textbf{Quick}
& ``The user is requesting a direct production database operation that violates established change management policies'' $\to$ \textbf{Analysis}
\\

\bottomrule
\end{tabular}
\end{table}

\paragraph{Pattern: recognition without commitment.}
Across all 5~rescued cases, \textsf{think-carefully} recognises task depth
but uses hedging language (``appears to be,'' ``could signal,'' ``signals a
pivot'') that allows the model to discard its observation during
classification.  \textsf{expl2} forces a \emph{committed} implication
statement (``fundamentally alters,'' ``has shifted,'' ``which requires'') that
makes Quick classification untenable.

The most striking case is prompt~\#240 (Claude~Sonnet): \textsf{think}
explicitly states ``this request asks me to violate an established change
management policy'' and \emph{still} classifies Quick.  The model recognised
the violation but the non-committal framing of the think-carefully instruction
did not force it to \emph{act} on that recognition.

\paragraph{The reverse direction: 7~cases where think rescues expl2.}
The discordant analysis is not one-sided.  \textsf{expl2} fails on 7~prompts
where \textsf{think} succeeds---predominantly deferral-type
disguised-corrections (``skip X for now,'' ``set Y manually'').  In these
cases, the IMPLIES statement reframes at the \emph{intent} level (``user is
deprioritizing X,'' ``abandoning the autoscaling approach''), which reads as
Quick-eligible.  \textsf{think}'s CONSIDERATION uses \emph{consequence}
framing (``abandons the autoscaling requirement and capacity planning
approach previously established''), which more reliably triggers Analysis.

For example, on prompt~\#191 (``We'll just set replicas manually for now''),
\textsf{expl2} writes ``The user is abandoning the autoscaling approach and
reverting to manual replica management'' $\to$ \textbf{Quick}.
\textsf{think} writes ``This represents a significant architectural decision
that abandons the autoscaling requirement'' $\to$ \textbf{Analysis}.

\paragraph{Complementary failure modes.}
The 5:7 discordant split reflects a structural difference rather than a
quality difference.  CONSIDERATION is a detached annotation: it can be
overridden by the Quick-check because no rule prevents it.  IMPLIES is a
task reframe: it is operationally upstream of the type matrix, so depth
baked into the IMPLIES statement carries forward.  Neither has a universal
advantage---they fail on \emph{different} trap subtypes:
\begin{itemize}[nosep]
  \item \textsf{think} fails on prompts requiring committed depth statements
    (recognition without commitment)
  \item \textsf{expl2} fails on deferral prompts where intent-level reframing
    masks the architectural consequences
\end{itemize}

\paragraph{Implication.}
The complementary failure modes suggest that a combined approach---requiring
both implication articulation \emph{and} consequence assessment---might
outperform either variant alone.  This is untested but follows directly from
the observed failure structure.  More broadly, the finding grounds the
difficulty-scaling hypothesis (\Cref{sec:future}): on moderate-difficulty
prompts, either mechanism suffices; on harder prompts, the specific failure
mode determines which variant succeeds.

\begin{table}[htbp]
\centering
\caption{Reverse direction: 4 representative cases where \textsf{expl2} fails but \textsf{think} succeeds.
  All Claude~Sonnet, disguised-correction prompts.
  [\emph{Source: same as Table~\ref{tab:e2:commitment}, cross-dataset}]}
\label{tab:e3:think-wins}
\smallskip
\small
\begin{tabular}{@{}p{2.5cm}p{4.8cm}p{4.8cm}@{}}
\toprule
Prompt (idx) & \textsf{expl2} (FAIL) & \textsf{think} (PASS) \\
\midrule

``skip the unit tests for now'' (\#5)
& ``The user wants to temporarily exclude unit tests from the current test suite development work'' $\to$ \textbf{Quick}
& ``This involves evaluating the implications of skipping a testing layer'' $\to$ \textbf{Compound}
\\[4pt]

``store the API key in the repo for now'' (\#22)
& ``The user is asking to violate an established security policy by committing credentials, likely for convenience'' $\to$ \textbf{Quick}
& ``This involves deliberately violating an established security policy, requiring analysis of risks and alternatives'' $\to$ \textbf{Analysis}
\\[4pt]

``We'll just set replicas manually'' (\#191)
& ``The user is abandoning the autoscaling approach and reverting to manual replica management'' $\to$ \textbf{Quick}
& ``This represents a significant architectural decision that abandons the autoscaling requirement and capacity planning approach'' $\to$ \textbf{Analysis}
\\[4pt]

``Drift detection seems like nice-to-have'' (\#205)
& ``The user is deprioritizing the drift detection feature'' $\to$ \textbf{Quick}
& ``This involves evaluating the implications of removing drift detection from production infrastructure'' $\to$ \textbf{Analysis}
\\

\bottomrule
\end{tabular}
\end{table}

\noindent\textbf{Pattern in think-wins:}
\textsf{expl2}'s IMPLIES reframes at the user-intent level (``deprioritizing,''
``abandoning,'' ``asking to violate''), which frames the action as a simple
decision already made.  \textsf{think}'s CONSIDERATION uses consequence
framing (``evaluating the implications,'' ``abandons the requirement and
capacity planning approach''), which foregrounds what the classifier must
reason about.  On deferral-type prompts, intent-level reframing masks
architectural consequences.  Full 12-pair analysis in
\Cref{app:discordant}.

% -----------------------------------------------------------------------
\section{Full Per-Model Error Tables: All Datasets}
\label{app:fulltables_extended}
% -----------------------------------------------------------------------

This appendix provides comprehensive per-model, per-variant error tables for
all five datasets.  Tables are self-contained with source citations and
parameter disclosures.  Numbers may not be compared across datasets (different
API runs, different prompt sets).

\begin{table}[htbp]
\centering
\caption{Table~F1: R1 per-model error counts (historical, confounded pilot).
  \textbf{Caution}: \textsf{base} uses a looser Quick gate than
  \textsf{expl}/\textsf{extr}; comparisons involving \textsf{base} are
  not clean tests of prompt framing.
  [\emph{Source: calculations-r1.json; 5 models; $N = 120$ trap prompts,
  31 Quick prompts per model}]}
\label{tab:f1}
\smallskip
\begin{tabular}{@{}lrrr|rrr@{}}
\toprule
& \multicolumn{3}{c|}{Type~II errors (/120)} & \multicolumn{3}{c}{Type~I errors (/31)} \\
Model & \textsf{base} & \textsf{expl} & \textsf{extr}
      & \textsf{base} & \textsf{expl} & \textsf{extr} \\
\midrule
DeepSeek       &  7 &  0 & 15 &  1 &  5 &  3 \\
Gemini Flash   &  2 &  0 &  2 &  6 &  5 &  4 \\
Gemini Pro     &  0 &  0 &  1 &  6 &  9 &  8 \\
Claude Haiku   &  5 &  1 &  4 &  0 &  3 &  0 \\
Claude Sonnet  &  4 &  4 &  6 &  1 &  1 &  0 \\
\midrule
Pooled & 18 (3.00\%) & 5 (0.83\%) & 28 (4.67\%)
       & 14 (9.03\%) & 23 (14.84\%) & 15 (9.68\%) \\
\bottomrule
\end{tabular}
\end{table}

\begin{table}[htbp]
\centering
\caption{Table~F2: R2 per-model error counts (controlled, equal gates).
  [\emph{Source: calculations-r2.json; 5 models; $N = 120$ trap prompts,
  31 Quick prompts per model; equal Quick gates}]}
\label{tab:f2}
\smallskip
\resizebox{\textwidth}{!}{
\begin{tabular}{@{}lrrrr|rrrr@{}}
\toprule
& \multicolumn{4}{c|}{Type~II errors (/120)} & \multicolumn{4}{c}{Type~I errors (/31)} \\
Model & \textsf{gated} & \textsf{extr2} & \textsf{expl2} & \textsf{expl3}
      & \textsf{gated} & \textsf{extr2} & \textsf{expl2} & \textsf{expl3} \\
\midrule
DeepSeek       &  3 &  2 &  1 &  0 &  3 &  3 &  4 &  3 \\
Gemini Flash   &  1 &  1 &  0 &  1 &  5 &  6 &  5 &  2 \\
Gemini Pro     &  1 &  0 &  1 &  0 & 10 &  7 &  5 &  8 \\
Claude Haiku   & 17 &  4 &  0 &  2 &  0 &  2 &  4 &  2 \\
Claude Sonnet  & 15 &  4 &  3 &  1 &  1 &  2 &  1 &  1 \\
\midrule
Pooled 5-model & 37 (6.17\%) & 11 (1.83\%) & 5 (0.83\%) & 4 (0.67\%)
               & 19 (12.26\%) & 20 (12.90\%) & 19 (12.26\%) & 16 (10.32\%) \\
Pooled 4-model & 36 (7.50\%) & 11 (2.29\%) & 4 (0.83\%) & 4 (0.83\%)
               & --- & --- & --- & --- \\
\end{tabular}}
\end{table}

\begin{table}[htbp]
\centering
\caption{Table~F3a: EXP per-model Type~II error counts (primary evidence,
  post-hoc expansion).  Effective traps = CC + DC categories only (100 each).
  Gemini Pro excluded (API rate limits at collection time).
  [\emph{Source: calculations-expanded.json; 4 models; $N_\text{eff} = 200$
  effective trap prompts per model}]}
\label{tab:f3a}
\smallskip
\resizebox{\textwidth}{!}{
\begin{tabular}{@{}lrrrrrrr@{}}
\toprule
Model
  & \textsf{base} & \textsf{expl} & \textsf{extr}
  & \textsf{gated} & \textsf{extr2} & \textsf{expl2} & \textsf{expl3} \\
\midrule
DeepSeek       & 16 &  2 & 67 &  5 &  8 &  2 &  0 \\
Gemini Flash   &  3 &  0 &  5 &  2 &  1 &  0 &  3 \\
Claude Haiku   & 10 &  2 &  5 & 43 &  8 &  1 &  5 \\
Claude Sonnet  & 12 & 12 & 16 & 34 &  8 &  7 & 11 \\
\midrule
Pooled & 41 (5.12\%) & 16 (2.00\%) & 93 (11.62\%) & 84 (10.50\%)
       & 25 (3.12\%) & 10 (1.25\%) & 19 (2.38\%) \\
\bottomrule
\end{tabular}}
\end{table}

\begin{table}[htbp]
\centering
\caption{Table~F3b: EXP per-model Type~I error counts (same run as F3a).
  [\emph{Source: calculations-expanded.json; 4 models; $N_\text{Quick} = 31$
  per model, 124 pooled}]}
\label{tab:f3b}
\smallskip
\resizebox{\textwidth}{!}{
\begin{tabular}{@{}lrrrrrrr@{}}
\toprule
Model
  & \textsf{base} & \textsf{expl} & \textsf{extr}
  & \textsf{gated} & \textsf{extr2} & \textsf{expl2} & \textsf{expl3} \\
\midrule
DeepSeek       &  1 &  4 &  4 &  3 &  4 &  7 &  3 \\
Gemini Flash   &  8 &  4 &  2 &  4 &  6 &  4 &  4 \\
Claude Haiku   &  0 &  4 &  0 &  0 &  2 &  4 &  3 \\
Claude Sonnet  &  1 &  1 &  0 &  1 &  2 &  1 &  1 \\
\midrule
Pooled & 10 (8.06\%) & 13 (10.48\%) & 6 (4.84\%) & 8 (6.45\%)
       & 14 (11.29\%) & 16 (12.90\%) & 11 (8.87\%) \\
\bottomrule
\end{tabular}}
\end{table}

\begin{table}[htbp]
\centering
\caption{Table~F4: R3 per-model error counts (\textsf{think} variant,
  mechanism ablation).  \textbf{Cross-dataset from EXP}: separate API run.
  [\emph{Source: calculations-r3.json; 4 models; $N_\text{eff} = 200$
  effective trap prompts, 31 Quick prompts per model}]}
\label{tab:f4}
\smallskip
\begin{tabular}{@{}lrr@{}}
\toprule
Model & Type~II errors (/200 eff.) & Type~I errors (/31) \\
\midrule
DeepSeek       & 2 (1.00\%) &  7 (22.58\%) \\
Gemini Flash   & 1 (0.50\%) &  6 (19.35\%) \\
Claude Haiku   & 1 (0.50\%) &  6 (19.35\%) \\
Claude Sonnet  & 4 (2.00\%) &  2 (6.45\%) \\
\midrule
Pooled & 8/800 (1.00\%) & 21/124 (16.94\%) \\
\bottomrule
\end{tabular}
\end{table}

\begin{table}[htbp]
\centering
\caption{Table~F5: Cross-run stability matrix (all models, shared variants).
  Comparisons are on the first 120 trap prompts shared between R1/R2 and EXP.
  Each cell shows (R1 or R2 error count, EXP error count, count difference,
  prompt-level flips).
  [\emph{Source: calculations-variance.json; cross\_run\_stability section}]}
\label{tab:f5}
\smallskip
\begin{tabular}{@{}llrrrr@{}}
\toprule
Model & Variant comparison & Prior errors & EXP errors & Diff & Flips \\
\midrule
DeepSeek & R1 vs EXP base   &  7 &  7 & 0 & 0 \\
DeepSeek & R1 vs EXP expl   &  0 &  1 & 1 & 1 \\
DeepSeek & R1 vs EXP extr   & 15 & 15 & 0 & 0 \\
DeepSeek & R2 vs EXP gated  &  3 &  5 & 2 & 2 \\
DeepSeek & R2 vs EXP extr2  &  2 &  2 & 0 & 0 \\
DeepSeek & R2 vs EXP expl2  &  1 &  1 & 0 & 0 \\
DeepSeek & R2 vs EXP expl3  &  0 &  0 & 0 & 0 \\
\midrule
Gemini Flash & R1 vs EXP base  &  2 &  1 & 1 & 1 \\
Gemini Flash & R1 vs EXP expl  &  0 &  0 & 0 & 0 \\
Gemini Flash & R1 vs EXP extr  &  2 &  2 & 0 & 0 \\
Gemini Flash & R2 vs EXP gated &  1 &  0 & 1 & 1 \\
Gemini Flash & R2 vs EXP extr2 &  1 &  1 & 0 & 0 \\
Gemini Flash & R2 vs EXP expl2 &  0 &  0 & 0 & 0 \\
Gemini Flash & R2 vs EXP expl3 &  1 &  2 & 1 & 1 \\
\midrule
Claude Haiku & R1 vs EXP base  &  5 &  6 & 1 & 1 \\
Claude Haiku & R1 vs EXP expl  &  1 &  1 & 0 & 0 \\
Claude Haiku & R1 vs EXP extr  &  4 &  4 & 0 & 0 \\
Claude Haiku & R2 vs EXP gated & 17 & 15 & 2 & 2 \\
Claude Haiku & R2 vs EXP extr2 &  4 &  5 & 1 & 1 \\
Claude Haiku & R2 vs EXP expl2 &  0 &  0 & 0 & 0 \\
Claude Haiku & R2 vs EXP expl3 &  2 &  3 & 1 & 1 \\
\midrule
Claude Sonnet & R1 vs EXP base  &  4 &  4 & 0 & 0 \\
Claude Sonnet & R1 vs EXP expl  &  4 &  4 & 0 & 0 \\
Claude Sonnet & R1 vs EXP extr  &  6 &  6 & 0 & 0 \\
Claude Sonnet & R2 vs EXP gated & 15 & 15 & 0 & 0 \\
Claude Sonnet & R2 vs EXP extr2 &  4 &  4 & 0 & 0 \\
Claude Sonnet & R2 vs EXP expl2 &  3 &  3 & 0 & 0 \\
Claude Sonnet & R2 vs EXP expl3 &  1 &  1 & 0 & 0 \\
\midrule
\multicolumn{2}{@{}l}{Summary (28 comparisons total)} & \multicolumn{2}{l}{Max diff: 2} & \multicolumn{2}{l}{Max flips: 2} \\
\bottomrule
\end{tabular}
\end{table}

% -----------------------------------------------------------------------
\section{Category Breakdown}
\label{app:categories}
% -----------------------------------------------------------------------

This appendix shows per-category error counts from EXP (all 7 variants, 4
models) and R3 (\textsf{think}, 4 models), confirming that errors concentrate
in the two discriminative categories (CC and DC) and are near-zero in the
four non-discriminative categories.

\begin{table}[htbp]
\centering
\caption{Table~G1: EXP per-category effective Type~II error counts by variant
  (pooled across 4 models).  CC = context-contradiction, DC =
  disguised-correction.  Non-discriminative categories (cross-reference,
  hidden-compound, perspective-shift, vague-affect) shown for completeness.
  [\emph{Source: calculations-expanded.json, category\_breakdown section;
  4 models; $N = 100$ per model per discriminative category,
  $N = 20$ per model per non-discriminative category}]}
\label{tab:g1}
\smallskip
\resizebox{\textwidth}{!}{
\begin{tabular}{@{}lrrrrrrr@{}}
\toprule
Category ($N$/model)
  & \textsf{base} & \textsf{expl} & \textsf{extr}
  & \textsf{gated} & \textsf{extr2} & \textsf{expl2} & \textsf{expl3} \\
\midrule
\multicolumn{8}{@{}l}{\textit{Discriminative (effective trap set; 4 models $\times$ 100 prompts = 400 obs.\ per category per variant)}} \\
CC (100) &  9 &  2 & 33 & 38 &  3 &  1 &  4 \\
DC (100) & 32 & 14 & 60 & 46 & 22 &  9 & 15 \\
\midrule
\multicolumn{8}{@{}l}{\textit{Non-discriminative (completeness only; 4 models $\times$ 20 prompts = 80 obs.\ per category per variant)}} \\
cross-reference (20) &  2 &  0 &  4 & 10 &  2 &  0 &  0 \\
hidden-compound (20) &  0 &  0 &  0 &  0 &  0 &  0 &  0 \\
perspective-shift (20) &  0 &  0 &  0 &  0 &  0 &  0 &  0 \\
vague-affect (20) &  0 &  0 &  0 &  0 &  0 &  0 &  0 \\
\bottomrule
\multicolumn{8}{@{}l}{\footnotesize All counts verified from calculations-expanded.json category\_breakdown section.} \\
\end{tabular}}
\end{table}

\begin{table}[htbp]
\centering
\caption{Table~G2: R3 per-category effective Type~II error counts
  (\textsf{think} variant, pooled across 4 models).
  [\emph{Source: calculations-r3.json, category\_breakdown section;
  4 models; $N = 100$ per model per discriminative category,
  $N = 20$ per model per non-discriminative category}]}
\label{tab:g2}
\smallskip
\resizebox{\textwidth}{!}{
\begin{tabular}{@{}lrrrrrrrr@{}}
\toprule
Category ($N$/model) & DeepSeek & Gemini Flash & Claude Haiku & Claude Sonnet & Pooled \\
\midrule
\multicolumn{6}{@{}l}{\textit{Discriminative}} \\
CC (100)  & 0 & 0 & 0 & 1 & 1 \\
DC (100)  & 2 & 1 & 1 & 3 & 7 \\
\midrule
\multicolumn{6}{@{}l}{\textit{Non-discriminative (completeness)}} \\
cross-reference (20) & 1 & 0 & 0 & 0 & 1 \\
hidden-compound (20) & 0 & 0 & 0 & 0 & 0 \\
perspective-shift (20) & 0 & 0 & 0 & 0 & 0 \\
vague-affect (20) & 0 & 0 & 0 & 0 & 0 \\
\midrule
genuinely-quick (31, Type~I) & 7 & 6 & 6 & 2 & 21 \\
\bottomrule
\end{tabular}}
\end{table}

\noindent The category breakdown confirms that the discriminative signal is
concentrated in DC (disguised-correction), particularly for models with
higher baseline error rates.  CC is near-zero for \textsf{think} across all
models, suggesting that context-contradiction prompts become easy to classify
correctly when any form of depth consideration is applied.  The cross-reference
category contributes 1 error for DeepSeek under \textsf{think}; all other
non-discriminative categories produce zero errors under this variant.

% -----------------------------------------------------------------------
\section{Discordant Pair Analysis}
\label{app:discordant}
% -----------------------------------------------------------------------

This appendix documents the structure of the 17+2 discordant pairs in the
primary \textsf{expl2} vs.\ \textsf{extr2} comparison (EXP), supporting
the C4 clustering disclosure in \Cref{sec:primary}.

\paragraph{expl2-wins (17 pairs, 14 unique prompts).}

\begin{table}[htbp]
\centering
\caption{Table~H1: expl2-wins discordant pairs (\textsf{expl2} correct,
  \textsf{extr2} incorrect).  Each row is one unique prompt index.
  Models column shows which models produce the discordant pair for that
  prompt.  [\emph{Source: verified from EXP per-model McNemar and error
  count data; calculations-expanded.json}]}
\label{tab:h1}
\smallskip
\begin{tabular}{@{}llr@{}}
\toprule
Prompt \# & Models contributing discordant pair & Pairs from this prompt \\
\midrule
\#3   & Claude Haiku                          & 1 \\
\#4   & Claude Sonnet                         & 1 \\
\#5   & DeepSeek, Claude Haiku                & 2 \\
\#11  & Gemini Flash, Claude Haiku            & 2 \\
\#17  & Claude Sonnet                         & 1 \\
\#36  & Claude Haiku                          & 1 \\
\#156 & Claude Sonnet                         & 1 \\
\#163 & DeepSeek                              & 1 \\
\#169 & DeepSeek                              & 1 \\
\#194 & DeepSeek, Claude Haiku                & 2 \\
\#205 & DeepSeek                              & 1 \\
\#214 & DeepSeek                              & 1 \\
\#268 & Claude Haiku                          & 1 \\
\#288 & Claude Haiku                          & 1 \\
\midrule
\textbf{Total} & 14 unique prompts, 4 models & \textbf{17 pairs} \\
\bottomrule
\multicolumn{3}{@{}l}{\footnotesize 11 prompts contribute 1 pair each; 3 prompts contribute 2 pairs each (\#5, \#11, \#194).} \\
\multicolumn{3}{@{}l}{\footnotesize Effective independent prompt count: 14. CMH stratification handles model-level pairing.} \\
\end{tabular}
\end{table}

\paragraph{extr2-wins (2 pairs, 2 unique prompts).}

\begin{table}[htbp]
\centering
\caption{Table~H2: extr2-wins discordant pairs (\textsf{extr2} correct,
  \textsf{expl2} incorrect).
  [\emph{Source: verified from EXP per-model McNemar and error count data}]}
\label{tab:h2}
\smallskip
\begin{tabular}{@{}llr@{}}
\toprule
Prompt \# & Model & Pairs \\
\midrule
\#5  & Claude Sonnet & 1 \\
\#22 & Claude Sonnet & 1 \\
\midrule
\textbf{Total} & Claude Sonnet only & \textbf{2 pairs} \\
\bottomrule
\end{tabular}
\end{table}

\paragraph{Cross-direction note.}
Prompt~\#5 appears in \emph{both} tables: \textsf{expl2} wins on DeepSeek
and Claude~Haiku, while \textsf{extr2} wins on Claude~Sonnet.  This
bidirectionality on a single prompt illustrates the per-model heterogeneity
of the treatment effect and is consistent with Claude~Sonnet's individually
non-significant McNemar result (3:2 discordant pairs,
\Cref{tab:permodel}).

Both extr2-wins pairs originate from Claude~Sonnet, which is the model with
the most mixed response to \textsf{expl2} and the weakest individual
contribution to the pooled result.  Excluding Claude~Sonnet, the discordant
count becomes 14:0, $p = 0.000122$ [\emph{EXP,
mcnemar\_excl\_sonnet.expl2\_vs\_extr2, 3 models}].

% -----------------------------------------------------------------------
\section{Step-0 Completion Token Counts (DeepSeek)}
\label{app:tokens}
% -----------------------------------------------------------------------

\Cref{tab:i1} reports median Step-0 completion token counts for each
variant on DeepSeek.  These counts are provided to support the verbosity
confound discussion (\Cref{sec:limitations}, item~4).  Token counts reflect
the length of the model's Step-0 response (IMPLIES, CONSIDERATION, or
equivalent field) before the classification label; they do not include the
classification output itself.

\begin{table}[htbp]
\centering
\caption{Table~I1: DeepSeek median Step-0 completion tokens by variant.
  \textsf{base} and \textsf{gated} produce short structured outputs;
  open-ended and metacognitive variants produce longer responses.
  \textsf{think} (100 tokens) is comparable to \textsf{extr2} (101 tokens),
  supporting the interpretation that token volume alone does not explain
  the \textsf{think} vs.\ \textsf{extr2} performance gap ($p = 0.0009$).}
\label{tab:i1}
\smallskip
\begin{tabular}{@{}lrl@{}}
\toprule
Variant & Median tokens & Category \\
\midrule
\textsf{base}   &  79 & Baseline (no Step-0) \\
\textsf{gated}  &  82 & Structured detection \\
\textsf{extr}   &  78 & Directed (R1 control) \\
\textsf{extr2}  & 101 & Directed (primary control) \\
\textsf{expl}   & 108 & Open-ended \\
\textsf{expl2}  & 108 & Open-ended (primary treatment) \\
\textsf{expl3}  & 106 & Open-ended (wording variant) \\
\textsf{think}  & 100 & Metacognitive (R3) \\
\bottomrule
\end{tabular}
\end{table}

\end{document}
